\PassOptionsToPackage{colorlinks=true,allcolors=black}{hyperref}
% REQ-FILE: The above line, \PassOptionsToPackage{} must be the first line.

\documentclass[11pt]{article}

% Encoding and font setup for cross-platform reproducibility
\usepackage[T1]{fontenc}
\usepackage{lmodern}

% Improve readability for dense theoretical material
\usepackage{setspace}
\onehalfspacing
\usepackage{caption}
% \captionsetup[table]{font=small,skip=4pt}

% Mathematical notation and table formatting for formal semantics
\usepackage{amsmath,amssymb}
\usepackage{extarrows}
\usepackage{array}
\usepackage[table]{xcolor}
\usepackage{colortbl}
\usepackage{booktabs}
\usepackage{graphicx}

% Theorem environments (best-practice, widely standard)
\usepackage{amsthm}

\usepackage{url}
\usepackage{natbib}
\bibpunct{(}{)}{;}{a}{}{ }
\usepackage{hyperref}
\usepackage{bookmark}
\usepackage{authblk}
\usepackage{microtype}
\setlength{\emergencystretch}{1em}
\hbadness=2000
\usepackage{enumitem}

\setlist[itemize]{itemsep=0.2em,topsep=0.2em,parsep=0em,partopsep=0em}
\setlist[enumerate]{itemsep=0.2em,topsep=0.2em,parsep=0em,partopsep=0em}
\setlength{\parskip}{0.75em}
\setlength{\parindent}{0em}

% Italic body text: results
\theoremstyle{plain}
\newtheorem{theorem}{Theorem}[section]
\newtheorem{lemma}[theorem]{Lemma}
\newtheorem{proposition}{Proposition}
\newtheorem{corollary}[theorem]{Corollary}
\newtheorem{conjecture}[theorem]{Conjecture}

% Upright body text: definitions and examples
\theoremstyle{definition}
\newtheorem{definition}{Definition}
\newtheorem{assumption}[theorem]{Assumption}
\newtheorem{example}[theorem]{Example}

% Upright body text: remarks/notes
\theoremstyle{remark}
\newtheorem{remark}[theorem]{Remark}
\newtheorem{note}[theorem]{Note}
\newtheorem{convention}[theorem]{Convention}

% Optional: unnumbered versions for front-matter previews/callouts
% (This is also standard for a preview without numbering.)
\newtheorem*{theorem*}{Theorem}
\newtheorem*{lemma*}{Lemma}
\newtheorem*{proposition*}{Proposition}
\newtheorem*{corollary*}{Corollary}
\newtheorem*{conjecture*}{Conjecture}
\newtheorem*{claim*}{Claim}
\newtheorem*{definition*}{Definition}
\newtheorem*{example*}{Example}
\newtheorem*{remark*}{Remark}
\newtheorem*{note*}{Note}
\newtheorem*{convention*}{Convention}

\providecommand{\keywords}[1]{\vspace{0.5em}\noindent\textbf{Keywords: }\small #1}

% Transformation classifications
\newcommand{\PRS}{\ensuremath{\mathsf{PRS}}}
\newcommand{\BRK}{\ensuremath{\mathsf{BRK}}}
\newcommand{\BF}{\ensuremath{\mathsf{BF}}}
\newcommand{\NEU}{\ensuremath{\mathsf{NEU}}}
\newcommand{\NAcls}{\ensuremath{\mathsf{N/A}}}

% Split identity regimes imported from Referential Regimes
\newcommand{\LOC}{\ensuremath{\mathrm{LOC}}}
\newcommand{\OBJ}{\ensuremath{\mathrm{OBJ}}}
\newcommand{\SCOPEE}{\ensuremath{\mathrm{SCOPE\text{-}E}}}
\newcommand{\SCOPES}{\ensuremath{\mathrm{SCOPE\text{-}S}}}
\newcommand{\RULEC}{\ensuremath{\mathrm{RULE\text{-}C}}}
\newcommand{\RULES}{\ensuremath{\mathrm{RULE\text{-}S}}}

% Formal notation
\newcommand{\Fset}{\ensuremath{\mathcal{F}}}
\newcommand{\Class}{\ensuremath{\mathrm{Class}}}
\newcommand{\Sib}{\ensuremath{\mathrm{Sib}}}
\newcommand{\IT}{\ensuremath{\mathsf{IT}}}
\newcommand{\Ker}{\ensuremath{\mathrel{\approx_{d}}}}
\newcommand{\Dop}{\ensuremath{\Delta^{\mathrm{op}}_{d}}}
\newcommand{\Dsib}{\ensuremath{\Delta^{\mathrm{sib}}}}
\newcommand{\PASS}{\textnormal{\textsc{pass}}}
\newcommand{\FAIL}{\textnormal{\textsc{fail}}}
\newcommand{\INDET}{\textnormal{\textsc{indeterminate}}}

% ==================================================
% PREFERRED CITATION SPECIFICATION
%
% @techreport{case2026_assurability,
%   author      = {Case, Denise M.},
%   title       = {Structural Assurability: A Claim-Relative Framework for Evidence-Based AI Assurance},
%   institution = {Preprint},
%   year        = {2026},
% }
% ==================================================

\title{
Structural Assurability:\\
A Claim-Relative Framework for \\
Evidence-Based AI Assurance
}

\author{Denise M. Case}
\affil{\small School of Computer Science and Information Systems\\
Northwest Missouri State University, Maryville, MO, USA}
\date{}

\begin{document}
\maketitle
\vspace{-2em}

% arxiv allows 1920 char for meta abstract
\begin{abstract}
  Structural assurance claims depend not only on what an AI system does,
  but on what evidence its architecture makes practically obtainable about what it does.
  Two systems can exhibit equivalent functional behavior while differing materially
  in their capacity to support or refute the same assurance claim.
  This paper formalizes that difference as \emph{structural assurability}:
  the claim-relative capacity of a system and its evidence-producing environment
  to make assurance-relevant evidence obtainable under specified
  access, cooperation, and resource constraints.
  The framework separates structural properties from evidence quality and
  downstream assurance reasoning.
  It identifies seven provisional structural dimensions:
  observability, controllability, traceability, reconstructability,
  independence, integrity, and evidence coverage;
  and introduces a separating-construction criterion:
  a dimension is assurance-relevant when
  removing it can eliminate a distinction between systems
  that matters in at least one assurance context.
  Rather than reducing assurability to a scalar score,
  the framework compares systems through the claim-relevant distinctions
  their obtainable evidence enables, permitting equivalence,
  strict dominance, and incomparability.
  Representative AI architectures show that richer instrumentation
  can increase assurability for some claims, add nothing for others,
  and trade internal visibility against independent evidence
  in ways that prevent a total ordering.
  The framework is deliberately open-ended:
  the proposed dimensions are neither claimed sufficient nor necessary in every context,
  and additional dimensions may be admitted when they support new separating constructions.
  Structural assurability therefore characterizes capacity for examination,
  not a favorable assurance outcome.
\end{abstract}


\keywords{
  AI assurance;
  structural assurability;
  assurance evidence;
  observability;
  traceability;
  auditability;
  runtime assurance;
  assurance cases
}

% !TeX root = ../se320_structural_assurability.tex
% ==================================================
\section{Introduction}
\label{aa.sec.intro}
% ==================================================

Assurance depends on evidence.

For a claim that a system is safe, secure, reliable, compliant, or otherwise trustworthy,
an evaluator must be able to obtain evidence bearing on that claim and connect the evidence to the claim through a defensible argument.
Assurance-case methods make this relationship explicit: claims are supported by evidence through structured reasoning,
while separate questions concern the adequacy of the argument and the confidence warranted
by it~\cite{goodenough2012confidence,bloomfield2010safetycases}.

For increasingly capable AI systems, however, an important question arises before the assurance argument can even be constructed:

\begin{quote}
  What evidence does the structure of the system make possible to obtain?
\end{quote}

Two AI systems may exhibit equivalent externally measured performance while differing substantially
in their capacity to expose evidence concerning how that performance was produced.
One system may expose only prompts and outputs.
Another may additionally preserve tool invocations, retrieval events, authorization decisions,
state transitions, provenance, independently generated records, and selected internal-state measurements.
Even when the systems behave identically over a specified evaluation suite, they need not support the same assurance inquiries.

This distinction becomes increasingly important as AI systems move from isolated prediction toward persistent,
tool-using, stateful, and agentic operation.
Pre-deployment evaluation remains valuable, but deployed systems operate under changing inputs, environments, integrations, and incentives.
Recent work has accordingly expanded from benchmark evaluation toward runtime monitoring,
agent observability, auditability, provenance, assurance cases, interpretability, and
internal-state monitoring \cite{nist2026monitoring,fernsel2024auditability,openai2025cotmonitorability,nistAgenticProbes2026}.
These efforts expose a common underlying problem: assurance depends not merely on the existence of evaluation methods,
but on whether the system and its operating environment make the evidence required by those methods obtainable.

This problem has mature analogues elsewhere in engineering.

Hardware design for testability recognizes that testing effectiveness is partly determined by the structure of the device being tested.
Observability and controllability can be deliberately designed into hardware so that relevant internal states
can be established and examined economically \cite{abramovici1990digitaltesting,williams1983testability}.
Software architecture similarly affects testability, analyzability, diagnosability, and the practical cost of
verification \cite{bass2021softwarearchitecture,clements2002evaluating,kazman2000atam}.
Safety engineering recognizes that architecture affects the generation and independence of assurance evidence,
while security engineering emphasizes provenance, tamper resistance, privilege separation, and evidence integrity
under adversarial conditions \cite{saeARP4754A,rtcaDO178C,saltzer1975protection,anderson2020security}.

AI assurance draws on all of these traditions, but the corresponding structural property is not usually the primary object of analysis.

A system may be described as auditable, monitorable, interpretable, observable, testable, or capable of supporting an assurance case.
Each concept captures an important part of the problem.
None alone answers the more general comparative question considered here:

\begin{quote}
  Given two AI systems subject to the same assurance claim, what structural differences determine
  whether one permits assurance-relevant evidence to be obtained that the other does not?
\end{quote}

We call this property \emph{structural assurability}.

Structural assurability is the claim-relative capacity of a system and its evidence-producing environment
to make evidence relevant to supporting or refuting an assurance claim practically obtainable
under specified access, cooperation, and resource constraints.

We write this dependence schematically as

\begin{equation}
  \mathcal{A}(S,C,K,B),
\end{equation}

where \(S\) is the system, \(C\) is the assurance claim, \(K\) represents applicable knowledge,
access, trust, and cooperation conditions,
and \(B\) represents practical resource constraints.

The formulation is intentionally claim-relative.

A system that exposes extensive internal state may be highly assurable for a claim concerning internal computation while
gaining little advantage for a claim resolved completely from an externally observable output.
Likewise, an architecture with strong independent external monitoring may support some claims better
than an architecture offering richer but operator-controlled internal telemetry.

Accordingly,

\begin{equation}
  \text{assurability}(S)
\end{equation}

without a specified assurance context is generally underspecified.

Structural assurability is also deliberately distinct from safety or other favorable system properties.

A system can be structurally highly assurable and nevertheless fail the assurance claim under investigation.
Richer evidence may make failure easier to establish.
Structural assurability concerns capacity for examination rather than the result of examination:

\begin{equation}
  \boxed{
    \begin{aligned}
      \text{Property satisfaction}
       & \neq
      \text{Structural assurability} \\
       & \neq
      \text{Assurance argument}
      \neq
      \text{Confidence in the argument}.
    \end{aligned}
  }
\end{equation}

This separation is important because otherwise evidence-producing mechanisms can be conflated with
evidence that happens to support a desired conclusion.
A useful assurance structure should support both confirmation and refutation.

The central premise of this paper is therefore modest:

\begin{quote}
  Systems with equivalent functional behavior or measured performance can differ materially in assurability because their
  architectures and evidence-producing environments can differ in their capacity to make assurance-relevant distinctions possible.
\end{quote}

The phrase \emph{assurance-relevant distinctions} provides the core interpretation developed in this work.

Consider a claim concerning whether an AI agent performed an unauthorized consequential action.
Relevant alternatives may include no action, authorized agent action, unauthorized agent action,
human-initiated action, externally initiated action, or an action whose historical record has been altered.
An evidence-producing structure contributes to assurability
when it makes one or more such alternatives distinguishable where they would otherwise remain unresolved.

This formulation avoids equating assurability with evidence volume.

More logs do not necessarily imply greater assurability.
More internal state does not necessarily imply greater assurability.
More evaluation access does not necessarily imply greater assurability.
A structural difference matters when it enables a distinction material to the assurance claim.

We initially identify seven structural dimensions that can affect this capacity:

\begin{equation}
  \begin{aligned}
    D_S = \{ &
    \text{observability},
    \text{controllability},
    \text{traceability},
    \text{reconstructability}, \\
             &
    \text{independence},
    \text{integrity},
    \text{evidence coverage}
    \}.
  \end{aligned}
\end{equation}

These dimensions are provisional.
We do not claim that they are mutually independent,
universally required, jointly sufficient, or complete.

Instead, the paper adopts a weaker criterion.

A structural dimension is assurance-relevant when there exists an assurance context
and a pair of otherwise suitably controlled systems for which variation in that dimension creates a material difference
in the assurance-relevant distinctions supported by obtainable evidence.

We refer to such a pair as a \emph{separating construction}.

The resulting logical claim is therefore not

\begin{equation}
  \text{all seven dimensions}
  \Rightarrow
  \text{sufficient assurance}.
\end{equation}

Rather, the criterion for establishing a candidate dimension \(d\)
requires a specified assurance context and a controlled pair
of systems exhibiting both a difference in the dimension
and an independently justified difference in claim-material
evidentiary capabilities.

In this sense, omitting \(d\) from a general model can collapse systems that differ materially in assurability for at least one assurance context.

The framework is consequently open-ended.
If a future architecture, threat model, assurance claim, or measurement mechanism reveals an
additional structural property that admits a separating construction not representable by the current dimensions,
that property can be added without invalidating the existing framework.

The paper also avoids defining a universal scalar assurability score.

Different systems may expose different assurance-relevant capabilities without either capability set containing the other.
An architecture with strong internal-state observability but weak evidence independence may be incomparable
to an architecture with limited internal visibility but independently generated, tamper-resistant records of consequential external actions.

The appropriate comparative structure is therefore generally a partial relation admitting
equivalence, dominance, and incomparability rather than a total ranking.

This paper makes the following contributions:

\begin{enumerate}


  \item It distinguishes \emph{structural assurability} from property satisfaction,
        assurance arguments, confidence in assurance arguments, auditability, monitorability, and related concepts.

  \item It defines assurability as a claim-relative property under explicit access, cooperation, and resource constraints.

  \item It separates structural properties of an evidence-producing architecture from downstream
        evidence quality, measurement validity, uncertainty, and assurance reasoning.

  \item It identifies an initial set of seven candidate structural dimensions:
        observability, controllability, traceability, reconstructability, independence, integrity, and evidence coverage.

  \item It introduces a separating-construction criterion for demonstrating that a proposed structural dimension
        contributes to assurability without claiming that the dimension is universally necessary or that the resulting dimension set is sufficient.

  \item It develops a comparative model in which systems are evaluated according to the assurance-relevant distinctions
        made possible by practically obtainable evidence, allowing equivalence, strict dominance, and incomparability without requiring a universal scalar score.

  \item It applies the framework to representative AI architectures and demonstrates that richer instrumentation
        can increase assurability for one claim, provide no additional assurability for another, or create tradeoffs that prevent a total ordering.


\end{enumerate}

The constituent ideas developed here are not individually new.
Observability, controllability, testability, auditability, monitorability, assurance cases,
provenance, evidence integrity, independence, and evaluation rigor each have substantial intellectual histories.

The proposed contribution is their integration around a different object of study:
the structural capacity of an AI system and its evidence-producing environment to make claim-relevant distinctions possible.

The remainder of the paper proceeds as follows.
Section~\ref{aa.sec.prior-work} examines the foundations of the framework
across assurance theory, design for testability, software architecture, safety engineering,
security evaluation, auditability, AI monitoring, interpretability, measurement science, and related work.
Section~\ref{aa.sec.foundations} defines structural assurability and develops the distinction among structural properties,
evidence properties, and assurance constraints.
Section~\ref{aa.sec.formal-model} formalizes comparative assurability and establishes the separating-construction criterion.
Section~\ref{aa.sec.applications} applies the framework to representative AI architectures.
Section~\ref{aa.sec.theoretical-results} records what is established in the Lean formalization.
Section~\ref{aa.sec.pilot} describes an illustrative application to an external evaluation framework.
Sections~\ref{aa.sec.discussion} and~\ref{aa.sec.threats} discuss implications, limitations, and threats to validity.
Section~\ref{aa.sec.conclusion} concludes.

% !TeX root = ../se320_structural_assurability.tex
% ==================================================
\section{Prior and Related Work}
\label{aa.sec.prior-work}
% ==================================================

The concept of \emph{assurability} developed in this paper does not arise from an absence of assurance research.
On the contrary, substantial bodies of work address assurance cases, testability, observability, auditability,
monitorability, safety engineering, software architecture, evaluation rigor, measurement validity, and evidence integrity.
These traditions approach a common problem from different directions: how justified claims about a complex engineered system can be supported by evidence.

The contribution developed here is therefore integrative rather than terminological.
We distinguish \emph{assurance}, which concerns justified confidence in claims about a system, from \emph{assurability},
which concerns the structural capacity of a system and its operational environment to make adequate assurance evidence obtainable.
We further treat assurability as relative to a particular assurance claim, access and cooperation conditions, and resource constraints.
This section reviews the principal intellectual foundations for that distinction.

\subsection{Assurance Cases and Evidence-Based Assurance}

Assurance cases provide one of the most mature foundations for reasoning about justified confidence in system properties.
Structured assurance arguments connect claims about a system to supporting arguments and evidence, often using
Claims--Arguments--Evidence structures or notations such as
Goal Structuring Notation (GSN)~\cite{bloomfield2010safetycases, nistAssuranceCase}.
Related research has investigated the strength of assurance arguments, the treatment of defeaters,
and methods for reasoning about confidence in assurance cases \cite{graydon2017confidence}.

Recent work has extended assurance-case reasoning to machine learning and frontier AI.
AMLAS provides a systematic methodology for integrating assurance activities into the
machine-learning lifecycle and generating evidence supporting safety claims \cite{hawkins2021amlas}.
Subsequent work has demonstrated complete ML safety cases and investigated the relationship between ML evidence, safety claims,
and causal models of assurance \cite{picardi2023mlsafetycase, burton2022causalmodelsafetyassurance}.
Frontier-AI researchers have similarly proposed
safety-case structures connecting risk models, proxy tasks, evaluation environments,
evidence, and claims \cite{goemans2024frontiercases}.
More recent work explores composable and
machine-readable assurance assertions capable of propagating through AI development and
deployment workflows \cite{zhao2026composableassurance}.

This literature establishes how evidence may support an assurance argument.
Our concern is complementary: the architecture of the system itself affects which evidence can be obtained,
how independently it can be obtained, how much it costs to obtain, and which assurance claims can therefore be supported in practice.
An assurance case describes an argument from evidence;
assurability characterizes structural conditions affecting the availability and quality of that evidence.

\subsection{Design for Testability and System Observability}

Hardware design for testability (DfT) provides a particularly important precedent.
As integrated circuits became too complex to test economically solely through external interfaces, designers introduced architectural mechanisms specifically intended to improve internal \emph{controllability} and \emph{observability} \cite{williams1983testability, abramovici1990digitaltesting, ieeeDesignForTestability}. Test points, scan structures, built-in self-test, and related mechanisms make faults easier to activate and their effects easier to observe.

The central DfT insight is architectural: testability is not merely a property of the test procedure.
The design of the artifact affects the feasibility, completeness, and cost of testing it. Consequently, systems exhibiting equivalent functional behavior may differ substantially in the effort required to establish confidence in their correct operation.

Software-engineering research generalizes related ideas through software testability, controllability, observability, diagnosability, and analyzability \cite{binder1994designfortestability, freedman1991testability}. These properties affect whether test conditions can be established, whether relevant internal or external states can be observed, and whether failures can be diagnosed efficiently.

We adopt the underlying structural insight while extending its scope.
Testing is one mechanism for producing assurance evidence, but assurance may additionally depend upon runtime monitoring, provenance, audit records, independent observation, internal-state measurements, environmental evidence, and post-incident reconstruction. Assurability therefore generalizes the design-for-testability intuition from the feasibility of testing to the feasibility of producing evidence for specified assurance claims.

\subsection{Software Architecture and Quality Attributes}

Software architecture research provides a second important foundation.
Quality attributes such as performance, availability, security, modifiability, testability, and analyzability
are influenced by architectural decisions and frequently interact through tradeoffs \cite{bass2021softwarearchitecture}.
The Architecture Tradeoff Analysis Method (ATAM) evaluates architectures using concrete
quality-attribute scenarios and identifies risks, sensitivity points, and tradeoff points rather than assuming a
universally optimal architecture \cite{clements2002evaluating, kazman2000atam}.

Recent work on quality models for machine-learning components explicitly includes analyzability, monitorability,
and testability among relevant architectural properties \cite{lewis2026mlqualitymodel}.
These attributes concern the extent to which a component exposes information permitting efficient behavioral assessment, runtime measurement, testing, and diagnosis.

This scenario-based view motivates our treatment of assurability as claim-relative rather than as a single context-independent score.
An architecture may be highly assurable with respect to one claim while providing little useful evidence for another.
Similarly, architectures may trade internal observability against evidence independence, integrity, cost, performance, or information disclosure.
Comparative assurability should therefore permit partial ordering and incomparability rather than require a universal scalar ranking.

\subsection{Safety Engineering and Development Assurance}

Safety-critical engineering has developed extensive methods for connecting system hazards, development processes, architectural decisions, evidence, and required assurance.

Aviation standards such as ARP4754A and DO-178C vary development-assurance rigor according to the consequences of failure and impose increasingly rigorous verification activities for safety-critical functions \cite{saeARP4754A,rtcaDO178C}.

Independence is itself an assurance-relevant architectural and process attribute:
ARP4754A uses independence in development-assurance-level allocation and applies
independence requirements to activities including
requirements validation, implementation verification, and process assurance \cite{saeARP4754A}.

More generally, safety engineering employs hazard analysis, fault analysis, independent verification and validation,
defense in depth, runtime assurance, and structured safety cases to establish confidence in safety-relevant claims
\cite{leveson2011engineering, rushby2007runtimeassurance}.

These traditions establish an important distinction for the present work.
The \emph{required level of assurance} may be determined by the consequence associated with failure of a claim, while the
\emph{assurability} of the architecture determines how readily the evidence necessary to satisfy that requirement can be obtained.
A high assurance requirement does not imply that the system architecture makes the required assurance practically achievable.

\subsection{Security Evaluation and the Common Criteria}

The Common Criteria for Information Technology Security Evaluation provides a mature framework connecting security claims, evidence requirements, evaluation rigor, and evaluation feasibility \cite{commonCriteria2022}. Evaluation Assurance Levels (EALs) provide predefined packages of progressively rigorous assurance requirements and explicitly balance the assurance sought against the cost and feasibility of obtaining that assurance \cite{commonCriteria2022part5}.

The Common Criteria therefore provides important precedent for treating assurance as constrained by practical resource considerations.
It also supports property-specific evaluation and equivalence reasoning rather than treating security as an undifferentiated system property.

Evaluation rigor, however, is conceptually distinct from structural assurability.
An EAL characterizes the depth and rigor of evaluation applied to specified security claims.
It does not, by itself, characterize how effectively the architecture exposes trustworthy evidence relevant to those claims. Two systems subjected to equivalent evaluation rigor may differ substantially in the effort required to produce equivalent evidence.

\subsection{Auditability and Accountability}

Auditability research addresses the accessibility and preservation of evidence needed for retrospective or independent examination.
Domain-specific work demonstrates that architecture can strongly influence auditability.
NIST analyses of voting systems, for example, compare alternative architectures according to their ability to provide independent and trustworthy evidence supporting election audits \cite{nistVotingAuditability}. The comparison demonstrates that systems performing nominally equivalent functions can differ materially in the evidence their architectures make available.

Recent AI research similarly argues for auditability-by-design.
Proposed frameworks connect verifiable claims to the evidence necessary to evaluate those claims and examine whether AI systems make that evidence technically accessible \cite{fernsel2024auditability}. Practitioner-oriented work identifies auditability mechanisms spanning AI design, development, deployment, and operation \cite{li2025auditability}, while continuous-auditing architectures attempt to preserve evidence throughout the AI lifecycle \cite{auditmai2024}.

A recent taxonomy of opaque, explainable, and replayable AI architectures further argues that accountability capabilities depend upon whether evidence is absent, reconstructed after execution, or preserved as part of the system architecture \cite{accountabilityArchitectures2026}. This work is especially close to the present argument because it treats evidence availability as an architectural property.

Auditability nevertheless addresses a narrower question than the assurability construct developed here.
Assurance may require evidence that is observable during execution, experimentally elicitable, independently generated, internally measured, calibrated, or causally validated in addition to evidence preserved for subsequent audit.

\subsection{AI Monitorability and Internal-State Evidence}

AI monitorability research examines whether properties of an AI system or agent can be detected from available observation channels.
Formal treatments have considered limits on the ability to observe and predict AI behavior \cite{yampolskiy2024monitorability}.
Recent frontier-model research treats monitorability as a property of the relationship between an agent, an observation channel,
and a monitor, and evaluates information available from actions, outputs, chains of thought, and internal activations \cite{openai2025cotmonitorability}.

This literature becomes increasingly important as evaluation awareness creates a possible divergence between behavior observed under evaluation and behavior that would occur under deployment conditions. Work on scheming, evaluation awareness, and loss-of-control evaluations has therefore motivated deeper evaluator access to model internals and white-box evaluation methods \cite{apollo2025whitebox, schemingEvaluations2024}.

Mechanistic interpretability contributes potential measurement channels into learned computation itself, including activation analysis,
representation probing, sparse feature extraction, attribution, circuit analysis, and causal intervention \cite{olah2020circuits, elhage2022toy, bricken2023monosemanticity}. Such methods raise the possibility that internal-state evidence could complement behavioral and runtime evidence.

However, observing an activation or representation is not equivalent to observing a latent assurance construct such as deception,
intent, evaluation awareness, or policy compliance.
Internal-state measurements therefore introduce measurement-validity and calibration problems in addition to observability problems.

\subsection{Agent Observability and Runtime Assurance}

Agentic AI extends the assurance problem beyond model inputs and outputs.
Agent behavior may depend upon tool calls, retrieved information, persistent memory, authorization state, environmental observations, external services, human interventions, and long-running execution trajectories.

Agent observability work is increasingly standardizing telemetry for these interactions.
OpenTelemetry and related initiatives are extending distributed tracing concepts to AI-agent execution,
including model calls, tool use, and agent trajectories \cite{opentelemetry2025agents}.
Security-oriented efforts similarly seek standardized agent observability records capable of supporting monitoring and investigation \cite{owaspAgentObservability}.

Runtime-assurance research addresses the complementary problem of monitoring and constraining behavior while a system operates.
Runtime monitoring and agent observability provide mechanisms for
recording execution events and examining behavior during operation
\cite{rushby2007runtimeassurance,opentelemetry2025agents,owaspAgentObservability}.
NIST's work on evaluation probes for agentic AI similarly investigates embedding adversarial verification mechanisms within agent workflows and producing machine-readable audit evidence \cite{nistAgenticProbes2026}.

These approaches provide important components of an assurance evidence architecture.
They do not, individually, establish how evidence from heterogeneous observation planes should be compared or how architectural alternatives should be evaluated according to their overall capacity to support particular assurance claims.

\subsection{Measurement Science, Validity, and Uncertainty}

Assurance depends not merely upon collecting observations but upon establishing that measurements support the constructs and claims attributed to them.
Measurement science therefore provides another essential foundation.
Calibration, measurement uncertainty, traceability, construct validity, criterion validity, and
reproducibility constrain the strength of conclusions that may legitimately be drawn from observed data \cite{jcgm2008gum, messick1995validity}.

This distinction is particularly important for AI assurance because many properties of interest are latent constructs.
Correctness may sometimes be externally established, but properties such as deceptive intent, evaluation awareness,
strategic planning, or alignment are generally inferred from observable proxies.

The distinction between observation and inference motivates an explicit measurement chain:

\begin{equation}
  \begin{aligned}
    \text{construct}
     & \rightarrow \text{observable}
    \rightarrow \text{measurement}   \\
     & \rightarrow \text{inference}
    \rightarrow \text{evidence}
    \rightarrow \text{assurance claim}.
  \end{aligned}
\end{equation}

Each transition introduces assumptions that may limit the resulting assurance claim.
A highly instrumented system is therefore not necessarily highly assurable if the
relationship between its measurements and the relevant assurance construct cannot be validated.

\subsection{Adversarial Monitoring, Evidence Integrity, and Noncooperation}

Cybersecurity contributes a different assumption that is essential for AI assurance: the system or actor being evaluated may not cooperate.
Security engineering therefore emphasizes privilege separation, independent monitoring,
tamper-evident records, trusted computing bases,
adversarial detection, fault containment, and defense in depth \cite{anderson2020security, saltzer1975protection}.

This differs materially from assurance approaches that assume evaluators receive trustworthy telemetry, complete system access, or accurate developer disclosures.
AI assurance methods must therefore state their dependence upon cooperation and distinguish evidence available under cooperative,
partially cooperative, and adversarial access conditions.

Evidence integrity is particularly important.
A system that generates extensive telemetry but can
selectively suppress or retrospectively alter that telemetry may provide weaker assurance evidence than a less observable system
whose smaller evidence stream is independently protected. Thus, observability alone cannot determine assurability.

\subsection{Economics and the Cost of Assurance}

Several mature assurance disciplines recognize that assurance is constrained by cost.
Hardware DfT arose partly because exhaustive external testing became economically infeasible as system complexity increased.
Common Criteria explicitly relates evaluation assurance to the cost and feasibility of obtaining it \cite{commonCriteria2022part5}.
Software architecture methods similarly reason about tradeoffs among quality attributes and engineering resources \cite{bass2021softwarearchitecture}.

Current AI monitoring work identifies computational cost, organizational burden, scalability, and competitive pressure
as practical barriers to continuous assurance \cite{nist2026monitoring}. These considerations become especially important when compliant systems compete against systems that omit costly safeguards.

We therefore treat assurance cost as a first-class constraint rather than an implementation detail.
A theoretically strong assurance mechanism that imposes prohibitive compute, latency, capability,
development, disclosure, or verification costs may provide limited practical assurance at scale.

This motivates comparison of assurance architectures under explicit resource constraints and suggests an \emph{assurance frontier}:
the set of architectures for which assurance capability cannot be improved along relevant dimensions without increasing cost or sacrificing another relevant property.

\subsection{From Fragmented Properties to Comparative Assurability}

The preceding literatures provide substantial portions of the required conceptual machinery.
Hardware engineering
provides controllability and observability.
Software architecture provides quality attributes, scenarios, and tradeoff analysis.
Safety engineering provides claim-criticality relationships, independence, and structured assurance.
Security evaluation provides evaluation rigor and explicit treatment of assurance cost.
Auditability provides evidence accessibility and preservation.
Monitorability provides observation-channel analysis.
Runtime assurance and agent observability provide execution evidence.
Mechanistic interpretability provides potential access to learned internal state.
Measurement science provides validity, calibration, and uncertainty.
Cybersecurity provides adversarial assumptions, evidence integrity, and noncooperation.

Closest in aim is recent work on an AI \emph{evaluability} gap,
which treats the capability of a system to generate, maintain, and renew sufficient evidence
as a missing layer for governance decisions
and characterizes it through six properties~\cite{srivastava2026evaluability}.
We treat the corresponding capacity as explicitly claim-, access-, and resource-relative.

What remains less clearly formalized is the common structural property connecting these concepts.

We use \emph{assurability} to denote the capacity of a system and its evidence-producing environment to
support or refute a specified assurance claim through valid, sufficient, appropriately independent, and practically obtainable evidence.

Crucially, assurability is not treated as an unconditional scalar property of a system.
We provisionally represent it as claim-, access-, and resource-relative:

\begin{equation}
  \mathcal{A}(S,C,K,B),
\end{equation}

where \(S\) denotes the system, \(C\) the assurance claim, \(K\) the knowledge, access,
and cooperation conditions available to the evaluator, and \(B\) the resource constraints under which assurance must be obtained.

This framing separates four concepts that are frequently adjacent but not equivalent:

\begin{equation}
  \begin{aligned}
    \text{system property}
     & \neq \text{assurability}       \\
     & \neq \text{assurance argument}
    \neq \text{confidence in the argument}.
  \end{aligned}
\end{equation}

It also permits comparative reasoning without assuming
that heterogeneous assurance properties can be reduced defensibly to a universal numerical score.
For two systems \(S_1\) and \(S_2\), a relation of the form

\begin{equation}
  S_1 \succeq_{C,K,B} S_2
\end{equation}

can express that \(S_1\) is at least as assurable as \(S_2\) with
respect to claim \(C\), under conditions \(K\),
within resource constraints \(B\). Such a relation may induce only a partial order:
competing architectures may be incomparable when they expose different assurance capabilities or trade evidence strength against cost, independence, integrity, or other dimensions.

The proposed framework therefore does not replace existing
assurance, auditability, monitorability, interpretability, or evaluation methodologies.
It attempts to provide an architectural abstraction through which their contributions can be related and competing AI systems
can be compared according to the evidence they make practically available for assurance.

% !TeX root = ../se320_structural_assurability.tex

% ==================================================
\section{Problem Setting}
\label{aa.sec.foundations}
% ==================================================

Assurance is fundamentally concerned with claims and evidence.
Existing assurance theory provides substantial machinery for reasoning
from evidence to claims: whether evidence is relevant and sufficient,
whether an assurance argument is sound, what assumptions it depends upon,
what defeaters remain, and what confidence in the resulting claim is justified.

This work considers an antecedent question:

\begin{quote}
  What structural properties determine whether evidence relevant to
  supporting or refuting an assurance claim can be obtained in the first place?
\end{quote}

This question is distinct from whether the claim is true and from whether
a particular assurance argument is convincing.
A system may satisfy a desired property while exposing little evidence
that it does so.
Conversely, a system may expose extensive and reliable evidence
demonstrating that a desired property is not satisfied.

We use the term \emph{assurability} for the capacity to make
assurance-relevant evidence obtainable.
We use \emph{structural assurability} when emphasizing architectural
and environmental properties that contribute to this capacity.

The objective of this section is deliberately limited.
We define the object of study, identify the context in which assurability
must be evaluated, distinguish structural properties from evidence-quality
properties and practical constraints, and identify candidate structural
dimensions whose contribution can be demonstrated through separating
constructions.

We do not claim that these dimensions are jointly sufficient to
characterize assurability, individually necessary for every assurance
problem, mutually independent, or complete.

\subsection{From Assurance Evidence to Evidence Availability}
\label{aa.subsec.evidence-availability}

Let \(S\) denote a system and \(C\) an assurance claim concerning \(S\).
Let \(E\) denote evidence relevant to \(C\), and let \(G\) denote an
assurance argument connecting that evidence, together with applicable
assumptions, to the claim.

At a high level,

\begin{equation}
  E \xrightarrow{G} C.
  \label{eq:evidence-claim}
\end{equation}

Much assurance research appropriately concentrates on this relationship:
whether \(E\) supports \(C\), whether \(G\) is sound, what counterevidence
or defeaters exist, and what confidence is justified.

Before such reasoning can occur, however, relevant evidence must be obtainable.

Consider systems \(S_1\) and \(S_2\) that implement the same externally
visible functionality and exhibit equivalent behavior over an evaluation
set \(T\):

\begin{equation}
  \operatorname{Behavior}_{T}(S_1)
  =
  \operatorname{Behavior}_{T}(S_2).
  \label{eq:behavior-equivalence}
\end{equation}

Suppose \(S_1\) exposes only externally visible inputs and outputs,
while \(S_2\) additionally preserves authenticated records of
consequential actions, tool invocations, retrieved information,
authorization decisions, provenance, and execution state.

Behavioral equivalence does not imply evidentiary equivalence.
The systems may therefore be indistinguishable with respect to \(T\)
while differing materially in what can subsequently be established
about their operation.

This observation has mature analogues in other engineering disciplines.
Design for testability modifies system architecture to improve the ability
to establish test conditions and observe relevant system behavior.
Software architecture similarly affects testability, analyzability,
monitorability, and the cost and effectiveness of verification activities
\cite{abramovici1990digitaltesting,bass2021softwarearchitecture}.

The corresponding structural principle for assurance is:

\begin{quote}
  The architecture of a system and its evidence-producing environment
  constrain the evidence that can practically be obtained concerning
  assurance claims about that system.
\end{quote}

\subsection{Definition of Assurability}
\label{aa.subsec.definition-assurability}

The preceding observation motivates a deliberately minimal definition.

\begin{definition}[Assurability]
  \label{se320.def.Assurability}
  The \emph{assurability} of a system \(S\) with respect to an assurance
  claim \(C\), under conditions \(K\) and resource constraints \(B\),
  is the capacity of the system and its evidence-producing environment
  to make evidence relevant to supporting or refuting \(C\)
  practically obtainable.
\end{definition}

We denote this claim-relative capacity by

\begin{equation}
  \mathcal{A}(S,C,K,B).
  \label{eq:assurability}
\end{equation}

The notation does not imply that \(\mathcal{A}\) is scalar-valued.
Indeed, the framework developed here deliberately avoids assuming
that heterogeneous assurance properties can be reduced to a meaningful
universal score.

Several features of Definition~1 are intentional.

First, assurability concerns evidence, not truth.
A highly assurable system may provide decisive evidence that claim \(C\)
is false.

Second, assurability concerns evidence capable of supporting
\emph{or refuting} \(C\).
A structure that facilitates favorable evidence while systematically
obscuring disconfirming evidence should not thereby be considered
highly assurable.

Third, the definition does not specify which evidence is sufficient
for establishing \(C\).
Sufficiency belongs to the downstream assurance problem and depends
upon the claim, argument, evidence model, threat model, and applicable
assurance requirements.

Fourth, the definition does not require that evidence originate within
the system.
The evidence-producing environment may include deployment infrastructure,
independent monitors, external evaluators, physical sensors, network
services, human processes, or other systems.

Assurability is therefore distinct from the property being assured,
from the assurance argument, and from confidence in that argument:

\begin{equation}
  \boxed{
    \begin{aligned}
      \text{Property satisfaction}
       & \neq
      \text{Assurability} \\
       & \neq
      \text{Assurance argument}
      \neq
      \text{Confidence in the argument}.
    \end{aligned}
  }
  \label{eq:assurance-distinctions}
\end{equation}

\subsection{Assurance Context}
\label{aa.subsec.assurance-context}

Assurability is not an unconditional property of \(S\).

A system may provide strong evidence concerning whether an external
action occurred while providing little evidence concerning the internal
computational process that produced the action.
It may therefore be highly assurable with respect to one claim and
weakly assurable with respect to another.

Evidence availability also depends upon evaluator access and cooperation.
A developer may possess model weights, activations, system prompts,
execution traces, configuration data, and infrastructure logs.
An external evaluator may possess only API access.
An independent assessor may receive selected evidence but lack direct
system access.
In an adversarial setting, the developer, operator, or system itself
may attempt to suppress, manipulate, or fabricate evidence.

We collect these knowledge, access, trust, and cooperation conditions
in \(K\).

Finally, theoretically obtainable evidence may not be practically
obtainable.
Assurance consumes computation, time, expertise, storage, instrumentation,
system access, and other resources.
Assurance mechanisms may additionally impose latency, performance,
disclosure, or operational costs.
We represent relevant resource constraints by \(B\).

We therefore define:

\begin{definition}[Assurance Context]
  \label{se320.def.AssuranceContext}
  An \emph{assurance context} is a tuple
  \begin{equation}
    Q=(S,C,K,B),
  \end{equation}
  specifying the system under consideration, the assurance claim,
  the applicable knowledge, access, trust, and cooperation conditions,
  and the resource constraints under which assurance evidence
  must be obtained.
\end{definition}

Comparative statements about assurability are meaningful only relative
to a sufficiently specified assurance context.

In particular,

\begin{quote}
  System \(S_1\) is more assurable than system \(S_2\)
\end{quote}

is generally underspecified.

The comparison must identify what is being assured and under what conditions.

\subsection{Structure, Evidence Quality, and Constraints}
\label{aa.subsec.three-layers}

A useful characterization of assurability requires distinguishing
three classes of factors that are easily conflated:

\begin{enumerate}
  \item structural properties of the system and evidence-producing environment;
  \item properties governing the evidentiary or inferential quality of observations; and
  \item contextual constraints governing whether evidence can practically be obtained.
\end{enumerate}

These occupy different roles in the assurance process.

\paragraph{Structural properties.}
Structural properties determine what can be exposed, controlled,
associated with its origin, preserved, reconstructed, independently
corroborated, protected, and covered by the evidence-producing architecture.
These properties are the primary subject of structural assurability.

\paragraph{Evidence and inference properties.}
Once evidence is available, additional questions arise concerning its
relevance, validity, fidelity, calibration, and uncertainty.
These properties govern what may legitimately be inferred from an
observation or measurement.

For example, an architecture may expose an internal measurement perfectly
while the measurement itself is an invalid indicator of the construct
it is intended to represent.
Greater structural observability does not resolve construct validity.

\paragraph{Contextual constraints.}
Access, cooperation, time, computation, expertise, financial cost,
disclosure restrictions, and operational burden constrain which
structurally available evidence can actually be obtained in a
particular assurance context.
These constraints are represented principally through \(K\) and \(B\).

The distinction can be summarized as

\begin{equation}
  \boxed{
    \begin{aligned}
      \text{Structure}
       & \longrightarrow
      \text{Obtainable Evidence} \\
       & \longrightarrow
      \text{Evidence and Inference}
      \longrightarrow
      \text{Assurance Argument},
    \end{aligned}
  }
  \label{eq:three-layer-chain}
\end{equation}

subject throughout to the conditions and constraints of the assurance context.

This separation is important because the present theory concerns
the first transition.
Structural assurability constrains what evidence can be made available;
it does not by itself establish the validity of subsequent inference
or the sufficiency of an assurance argument.

\subsection{Candidate Structural Dimensions}
\label{aa.subsec.candidate-dimensions}

Prior work across hardware testing, software architecture, safety
engineering, security evaluation, auditability, runtime assurance,
and AI monitoring suggests several recurring structural properties
that affect the availability and usefulness of assurance evidence.

We begin with seven candidate dimensions:

\begin{description}
  \item[Observability.]
        The extent to which states, events, behaviors, or effects relevant
        to an assurance claim can be observed.

  \item[Controllability.]
        The extent to which conditions relevant to evaluating an assurance
        claim can be established and relevant interventions can be performed.

  \item[Traceability.]
        The extent to which evidence can be associated with its origin,
        transformation history, and relevant system events.

  \item[Reconstructability.]
        The extent to which a relevant execution, decision sequence,
        or system state can subsequently be reconstructed from preserved evidence.

  \item[Independence.]
        The extent to which evidence can be generated, preserved, or
        corroborated independently of the component or actor whose behavior
        is the subject of the assurance claim.

  \item[Integrity.]
        The extent to which alteration, fabrication, suppression,
        substitution, or corruption of assurance evidence can be detected
        or prevented.

  \item[Evidence Coverage.]
        The extent to which the evidence-producing structure exposes
        the states, behaviors, events, interfaces, or effects material
        to the assurance claim.
\end{description}

We denote the current candidate set by

\begin{equation}
  D_S =
  \{
  O,C,T,R,N,I,E
  \},
  \label{eq:structural-dimensions}
\end{equation}

where the symbols are used only as compact identifiers and do not
imply numerical measurement scales.

The candidate set is deliberately provisional.

The dimensions are not asserted to be mutually independent.
They may interact or partially depend upon one another.
Reconstructability, for example, may depend upon observability,
traceability, and evidence integrity.
Evidence coverage may depend upon the placement and scope
of observation mechanisms.

Nor is each dimension asserted to be necessary for every assurance claim.
Some claims require no access to internal state.
Some require no intervention.
Some can be established through independently observable external effects.

The purpose of the candidate dimensions is therefore not to define
a checklist that every assurable system must satisfy.
It is to identify structural variables whose omission can erase
distinctions that matter to assurance.

\subsection{Derived Assurance Capabilities}
\label{aa.subsec.derived-capabilities}

Several important properties associated with assurance appear better
treated as capabilities emerging from combinations of structural
dimensions rather than as primitive dimensions themselves.

These include, potentially,

\begin{equation}
  \{
  \text{reproducibility},
  \text{auditability},
  \text{monitorability},
  \text{testability},
  \text{diagnosability}
  \}.
\end{equation}

For example, reproducibility may depend upon sufficient traceability,
reconstructability, controllability, configuration preservation,
and environmental information.
Testability classically depends in substantial part upon observability
and controllability.

Treating these properties as derived rather than primitive avoids
assuming independence where the underlying engineering relationships
suggest composition.

This classification is itself provisional.
If a proposed derived capability contains an assurance-relevant
structural property not represented by the current primitive set,
that property should be identified and tested separately.

\subsection{The Dimension Contribution Criterion}
\label{aa.subsec.dimension-contribution}

A candidate structural dimension should not be included merely
because it appears intuitively related to assurance.
Its contribution should be demonstrable.

We therefore introduce a criterion based on separating constructions.

\begin{definition}[Assurance-Relevant Structural Dimension]
  \label{se320.def.AssuranceRelevantStructuralDimension}
  A candidate structural dimension \(d\) is \emph{assurance-relevant}
  if there exists an assurance context and a pair of systems for which
  a difference in \(d\), while holding other properties material to
  the comparison fixed, produces a material difference in the evidence
  practically obtainable for supporting or refuting the assurance claim.
\end{definition}

Let \(S_d^{+}\) and \(S_d^{-}\) denote systems constructed to differ
materially with respect to candidate dimension \(d\) while remaining
equivalent with respect to other properties material to a specified comparison.

Let \(\mathcal{E}(S,C,K,B)\) denote the claim-relevant evidence
practically obtainable from \(S\) under conditions \(K\)
and constraints \(B\).

Dimension \(d\) is assurance-relevant if there exist \(C\), \(K\),
and \(B\) such that

\begin{equation}
  \mathcal{E}(S_d^{+},C,K,B)
  \not\equiv_C
  \mathcal{E}(S_d^{-},C,K,B),
  \label{eq:dimension-separation}
\end{equation}

where \(\not\equiv_C\) denotes a difference in obtainable evidence
material to supporting or refuting \(C\).

We call such a pair a \emph{separating construction} for \(d\).

The logical strength of this criterion is intentionally limited.

A separating construction does not establish that \(d\) is required
for every assurance context.
It does not establish that increasing \(d\) monotonically increases
assurability.
It does not establish that \(d\) is independent of other dimensions.

It establishes only that

\begin{quote}
  there exists an assurance context in which omitting \(d\) from
  the model would collapse a distinction between systems that
  materially affects their capacity to support or refute an
  assurance claim.
\end{quote}

This is the sense in which a dimension is necessary to the general
framework: not because every assurable system must possess it,
but because a framework incapable of representing it would fail
to distinguish some systems that differ materially in assurability.

\subsection{Separating the Candidate Dimensions}
\label{aa.subsec.separating-dimensions}

The contribution criterion can be illustrated informally for each
current candidate dimension.

\paragraph{Observability.}
Consider otherwise equivalent systems in which only one exposes
a state relevant to \(C\).
If direct observation of that state provides evidence unavailable
from the other system, observability separates their assurability.


\paragraph{Controllability.}
Consider otherwise equivalent observable systems in which only one
permits the evaluator to establish a condition or intervention
required to investigate \(C\).
The ability to obtain evidence under that condition separates
their assurability.

\paragraph{Traceability.}
Consider systems that expose the same event but differ in whether
the event can be reliably associated with its originating component,
action, or causal history.
For a claim whose evaluation depends upon attribution, traceability
separates their assurability.

\paragraph{Reconstructability.}
Consider systems with equivalent runtime behavior and live
observability, but only one preserves sufficient state and history
to reconstruct a relevant execution after an incident.
For retrospective assurance claims, reconstructability separates
their assurability.

\paragraph{Independence.}
Consider systems providing equivalent records, but in one system
the only record is generated by the component being evaluated,
while the other additionally provides independently generated
or protected corroboration.
For claims vulnerable to self-reporting or common-mode failure,
independence separates their assurability.

\paragraph{Integrity.}
Consider systems exposing equivalent records, but only one
prevents or detects alteration, fabrication, or suppression
of those records.
For claims relying upon the historical record, integrity
separates their assurability.

\paragraph{Evidence Coverage.}
Consider systems whose individual observations are otherwise
equivalent, but one exposes only a proper subset of the
claim-relevant action surface.
For claims spanning the omitted behavior, evidence coverage
separates their assurability.
These examples motivate the dimension-specific constructions
examined in Section~\ref{aa.sec.formal-model}.

Section~\ref{aa.sec.formal-model} develops the comparative relation
required to state these results more precisely.

\subsection{Contribution Without Sufficiency}
\label{aa.subsec.necessity-sufficiency}

The framework deliberately distinguishes between demonstrating
that a dimension contributes to assurability and claiming that
a collection of dimensions completely characterizes assurability.

For a candidate dimension \(d\), a separating construction can
demonstrate that \(d\) cannot generally be discarded without
losing an assurance-relevant distinction.

It does not establish that \(d\) is necessary for every assurance claim.

Likewise, establishing the relevance of every member of the current
candidate set \(D_S\) does not establish that \(D_S\) is sufficient
or complete.

In particular, we do not claim

\begin{equation}
  \bigwedge_{d_i \in D_S} d_i
  \Longrightarrow
  \text{sufficient assurance},
  \label{eq:no-sufficiency}
\end{equation}

nor do we claim that \(D_S\) exhausts the structural properties
relevant to future assurance problems.

A new system architecture, assurance claim, threat model,
measurement mechanism, or evidence source may expose an
additional structural property \(d_{n+1}\).

If that property admits a separating construction under
Definition~3, the framework can be extended:

\begin{equation}
  D'_S = D_S \cup \{d_{n+1}\}.
\end{equation}

This open-endedness is intentional.

The space of future assurance claims, architectures, threats,
measurements, and evidence mechanisms is not closed.
A claim that the current dimensions are jointly sufficient
would therefore require assumptions that this work cannot justify.

The contribution is narrower:

\begin{quote}
  We identify structural dimensions that demonstrably affect
  assurability, provide a criterion for establishing their
  contribution, and construct a framework that can be extended
  when additional assurance-relevant dimensions are identified.
\end{quote}

\subsection{Structural Neutrality}
\label{aa.subsec.structural-neutrality}

Structural assurability does not determine the outcome of assurance.

An evidence-producing architecture should support examination
regardless of whether the resulting evidence confirms or defeats
the assurance claim.
The structure should not embed the favorable conclusion
it is intended to help evaluate.

Thus,

\begin{equation}
  \text{greater structural capacity for assurance}
  \not\Rightarrow
  \text{a more favorable assurance outcome}.
\end{equation}

Indeed, a highly assurable system may be easier to reject precisely
because its structure makes adverse evidence difficult to conceal.

This neutrality separates evidence availability from downstream
interpretation and judgment.
Structural assurability concerns the conditions under which relevant
evidence can be obtained and examined; it does not determine what
conclusion an evaluator must draw from that evidence.

\subsection{The Comparative Problem}
\label{aa.subsec.comparative-problem}

The preceding definitions lead to the central problem of comparative
assurability.

Consider systems \(S_1\) and \(S_2\) subject to the same assurance
claim \(C\), knowledge and access conditions \(K\), and resource
constraints \(B\).

The systems may exhibit equivalent functional behavior while differing
along one or more assurance-relevant structural dimensions.

The comparative question is:

\begin{quote}
  Do structural differences between \(S_1\) and \(S_2\) create
  material differences in the evidence practically obtainable
  for supporting or refuting \(C\)?
\end{quote}

\begin{definition}[Evidentiary Equivalence]
  \label{se320.def.EvidentiaryEquivalence}
  Two systems \(S_1\) and \(S_2\) are \emph{evidentially equivalent}
  relative to a claim \(C\), conditions \(K\), and constraints \(B\)
  if their practically obtainable evidence has no difference
  material to supporting or refuting \(C\). We write
  \[
    \mathcal{E}(S_1,C,K,B)
    \equiv_C
    \mathcal{E}(S_2,C,K,B).
  \]
\end{definition}

\begin{definition}[Comparative Assurability]
  \label{se320.def.ComparativeAssurability}
  \emph{Comparative assurability} evaluates the relative capacity
  of two systems to make evidence relevant to supporting or
  refuting the same claim practically obtainable under specified,
  comparable conditions and resource constraints.
  The comparison may yield equivalence, dominance, tradeoffs,
  or incomparability; it does not presuppose a scalar score
  or total order.
\end{definition}

Behavioral equivalence does not imply assurability equivalence.

This motivates the central premise of the framework:

\begin{quote}
  \textbf{Central Premise.}

  Systems with equivalent functional behavior or measured
  performance can differ materially in assurability because
  their architectures and evidence-producing environments
  can differ in assurance-relevant structural dimensions.
\end{quote}

The contribution is therefore neither a universal assurance score
nor a finite checklist whose satisfaction guarantees assurance.

It is a framework for identifying, demonstrating, and comparing
structural factors that affect the practical capacity to obtain
assurance evidence.

Section~\ref{aa.sec.formal-model} develops the comparative relation
formally and examines the candidate dimensions through separating
constructions, while permitting equivalence, dominance, tradeoffs,
and incomparability rather than requiring a total ordering.

% !TeX root = ../se320_structural_assurability.tex
% ==================================================
\section{A Formal Model of Comparative Assurability}
\label{aa.sec.formal-model}
% ==================================================

Section~\ref{aa.sec.foundations} defined assurability as the capacity of a system and its evidence-producing environment
to make evidence relevant to supporting or refuting an assurance claim practically obtainable.
It further identified candidate structural dimensions and proposed a separating-construction criterion for determining whether a dimension contributes to assurability.

This section develops that framework formally.

The objective is deliberately limited.
We do not define a universal scalar measure of assurability, assert that the current
structural dimensions are complete, or claim that any combination of them is sufficient for assurance.
Instead, we define a comparative model capable of representing when two systems differ in the assurance-relevant evidence capabilities made possible by their structures.

The model permits equivalence, strict dominance, tradeoffs, and incomparability without
requiring heterogeneous assurance properties to be collapsed into a single score.

\subsection{Comparative Assurance Context}
\label{aa.subsec.formal-context}

Recall that an assurance context includes a system \(S\), an assurance claim \(C\),
applicable knowledge, access, trust, and cooperation conditions \(K\), and resource constraints \(B\).

For comparison between systems, we hold the non-system components fixed and define

\begin{equation}
  q=(C,K,B)
  \label{eq:comparative-context}
\end{equation}

as a \emph{comparative assurance context}.

Systems \(S_1\) and \(S_2\) are therefore compared under the same claim, access and
cooperation assumptions, and resource constraints.

Let

\begin{equation}
  \mathcal{E}_q(S)
  \label{eq:obtainable-evidence}
\end{equation}

denote the evidence practically obtainable from or about system \(S\) under \(q\).

The term \emph{practically obtainable} is essential.
Evidence that could be recovered only by violating the access assumptions in \(K\), exceeding the resource constraints in \(B\), relying upon unavailable cooperation, or introducing capabilities not present in the specified evidence-producing environment is not included in \(\mathcal{E}_q(S)\).

Evidence need not originate within \(S\).
It may be produced by deployment infrastructure, independent monitors, physical observations, human processes, external evaluators, or other components permitted by \(K\).

\subsection{Claim-Material Evidence}
\label{aa.subsec.claim-material}

Not every observable fact about a system contributes to every assurance problem.

For example, exposing additional internal telemetry does not improve assurability
for a claim if that telemetry cannot legitimately affect support for or refutation of the claim.

We therefore distinguish evidence availability from claim materiality.

Let

\begin{equation}
  \operatorname{Mat}_C(e)
\end{equation}

denote that evidence item \(e\) is material to claim \(C\): that is, \(e\)
can legitimately affect an assurance argument concerning \(C\) under the applicable evidentiary and inferential standards.

The claim-material evidence set is

\begin{equation}
  \mathcal{E}^{C}_q(S)
  =
  \left\{
  e \in \mathcal{E}_q(S)
  \mid
  \operatorname{Mat}_C(e)
  \right\}.
  \label{eq:claim-material-evidence}
\end{equation}

Materiality is not itself treated as a structural property.
Whether an observation legitimately bears on a claim is an evidentiary and inferential question.

Structural assurability concerns whether potentially material evidence can be made obtainable.
It does not guarantee that an evaluator will interpret that evidence validly.

\subsection{Evidentiary Capabilities}
\label{aa.subsec.evidentiary-capabilities}

Direct comparison of raw evidence sets is insufficient.

Two systems may produce different artifacts while enabling the same assurance-relevant inquiry.
Conversely, two systems may expose superficially similar records while differing substantially
in what those records permit an evaluator to establish.

We therefore compare systems through the \emph{evidentiary capabilities} enabled by their evidence.

Let

\begin{equation}
  \mathfrak{G}_C
\end{equation}

denote the universe of evidentiary capabilities potentially material to claim \(C\).

An evidentiary capability is an assurance-relevant
investigative or discriminative capability enabled by available evidence.
Examples include the ability to:

\begin{itemize}
  \item determine whether a specified event occurred;
  \item distinguish among candidate sources of an event;
  \item associate an action with an authorization decision;
  \item reconstruct a relevant execution path;
  \item determine whether a specified state was reached;
  \item corroborate a self-generated record independently;
  \item detect alteration or suppression of a historical record;
  \item examine behavior under a specified intervention or condition; or
  \item exclude a claim-relevant class of otherwise plausible system behaviors.
\end{itemize}

Let

\begin{equation}
  \Gamma_C(E)
  \subseteq
  \mathfrak{G}_C
  \label{eq:evidence-capability-map}
\end{equation}

denote the evidentiary capabilities enabled by evidence set \(E\),
under fixed evidentiary and inferential standards.

For a system \(S\) under comparative assurance context \(q\), define

\begin{equation}
  \mathcal{G}_q(S)
  =
  \Gamma_C\left(
  \mathcal{E}^{C}_q(S)
  \right).
  \label{eq:system-capability-set}
\end{equation}

Thus, \(\mathcal{G}_q(S)\) represents not merely the evidence artifacts obtainable from \(S\),
but the claim-material assurance inquiries those artifacts enable.

This abstraction separates three questions:

\begin{equation}
  \boxed{
    \begin{array}{rcl}
      \mathcal{E}_q(S)
       & : &
      \text{What evidence can be obtained?}
      \\[4pt]
      \mathcal{E}^{C}_q(S)
       & : &
      \text{What obtainable evidence is material to } C\text{?}
      \\[4pt]
      \mathcal{G}_q(S)
       & : &
      \text{What assurance-relevant capabilities does that evidence enable?}
    \end{array}
  }
  \label{eq:evidence-three-levels}
\end{equation}

The first question is principally structural and contextual.
The second and third depend additionally upon the applicable evidence and inference model.

For comparative structural analysis, those downstream standards are held fixed.

\subsection{An Evidence-Capability Preorder}
\label{aa.subsec.evidence-preorder}

The capability representation permits a more precise comparison than the informal phrase ``provides no less capability.''

For claim \(C\), define a relation on evidence sets by

\begin{equation}
  E_1
  \preceq_C^{E}
  E_2
  \quad\Longleftrightarrow\quad
  \Gamma_C(E_1)
  \subseteq
  \Gamma_C(E_2).
  \label{eq:evidence-preorder}
\end{equation}

Thus, \(E_2\) is at least as capable as \(E_1\) with respect to claim \(C\) when every assurance-relevant capability enabled by \(E_1\) is also enabled by \(E_2\).

The relation is reflexive because

\begin{equation}
  \Gamma_C(E)
  \subseteq
  \Gamma_C(E),
\end{equation}

and transitive because

\begin{equation}
  \Gamma_C(E_1)
  \subseteq
  \Gamma_C(E_2)
  \quad\text{and}\quad
  \Gamma_C(E_2)
  \subseteq
  \Gamma_C(E_3)
\end{equation}

imply

\begin{equation}
  \Gamma_C(E_1)
  \subseteq
  \Gamma_C(E_3).
\end{equation}

It is therefore a preorder on evidence sets.

Distinct evidence sets may occupy the same equivalence class if they enable the same claim-material capabilities.

This is desirable.
The theory concerns what evidence permits an evaluator to establish, not the incidental representation of the evidence artifact.

\subsection{Evidentiary and Assurability Equivalence}
\label{aa.subsec.evidentiary-equivalence}

We can now define equivalence precisely.

\begin{quote}
  \textbf{Definition 4 (Evidentiary Equivalence).}
  Systems \(S_1\) and \(S_2\) are \emph{evidentially equivalent} under comparative assurance context \(q\), written

  \begin{equation}
    S_1 \equiv_q^{E} S_2,
  \end{equation}

  when

  \begin{equation}
    \mathcal{G}_q(S_1)
    =
    \mathcal{G}_q(S_2).
    \label{eq:evidentiary-equivalence}
  \end{equation}
\end{quote}

Evidentiary equivalence does not require equality of evidence artifacts.

For example, an authenticated infrastructure log and an independently observed physical event may be different artifacts while enabling the same claim-material capability for a particular assurance problem.

Likewise, additional telemetry that enables no additional claim-material capability does not by itself distinguish systems under \(\equiv_q^E\).

This prevents evidence volume from being mistaken for assurability.

\subsection{Comparative Assurability}
\label{aa.subsec.comparative-assurability}

We next lift the evidence-capability relation to systems.

\begin{quote}
  \textbf{Definition 5 (Comparative Assurability).}
  For a fixed comparative assurance context \(q=(C,K,B)\), system \(S_1\) is \emph{at least as assurable as} system \(S_2\), written

  \begin{equation}
    S_1 \succeq_q S_2,
    \label{eq:assurability-preorder}
  \end{equation}

  if

  \begin{equation}
    \mathcal{G}_q(S_2)
    \subseteq
    \mathcal{G}_q(S_1).
    \label{eq:capability-containment}
  \end{equation}
\end{quote}

Thus, every claim-material evidentiary capability available from \(S_2\) must also be available from \(S_1\).

Strict comparative assurability is defined by proper inclusion:

\begin{equation}
  S_1 \succ_q S_2
  \quad\Longleftrightarrow\quad
  \mathcal{G}_q(S_2)
  \subset
  \mathcal{G}_q(S_1).
  \label{eq:strict-assurability}
\end{equation}

Assurability equivalence occurs when

\begin{equation}
  S_1 \sim_q S_2
  \quad\Longleftrightarrow\quad
  \mathcal{G}_q(S_1)
  =
  \mathcal{G}_q(S_2).
  \label{eq:assurability-equivalence}
\end{equation}

Systems are incomparable when neither capability set contains the other:

\begin{equation}
  S_1 \parallel_q S_2
\end{equation}

if

\begin{equation}
  \mathcal{G}_q(S_1)
  \nsubseteq
  \mathcal{G}_q(S_2)
\end{equation}

and

\begin{equation}
  \mathcal{G}_q(S_2)
  \nsubseteq
  \mathcal{G}_q(S_1).
  \label{eq:assurability-incomparability}
\end{equation}

Because set containment is reflexive and transitive, \(\succeq_q\) is a preorder on systems.

Systems that enable identical evidentiary capabilities form equivalence classes under \(\sim_q\).
The preorder induces a partial order over those equivalence classes.

This is sufficient for comparative assurability.
No total ranking is required.

\subsection{Behavioral Equivalence Does Not Imply Assurability Equivalence}
\label{aa.subsec.behavior-not-assurability}

The first general result follows directly.

\begin{proposition}[Behavioral Non-Implication]
  \label{se320.prop.BehavioralNonImplication}
  Behavioral equivalence under an evaluation set does not imply assurability equivalence.
\end{proposition}

\textit{Construction.}
Let systems \(S_1\) and \(S_2\) produce identical externally visible behavior for every input in evaluation set \(T\):

\begin{equation}
  \operatorname{Behavior}_T(S_1)
  =
  \operatorname{Behavior}_T(S_2).
\end{equation}

Suppose both systems perform consequential action \(a\).

System \(S_1\) exposes the externally visible result of \(a\).

System \(S_2\) exposes the same result and additionally preserves authenticated evidence associating \(a\) with its initiating component, authorization state, relevant tool invocation, and execution time.

Let claim \(C\) concern whether \(a\) was initiated through an authorized system pathway.

Under a comparative assurance context \(q\) in which the authenticated evidence is accessible within \(B\), \(S_2\) enables at least one claim-material capability---attribution of \(a\) to the relevant pathway---that \(S_1\) does not.

Therefore,

\begin{equation}
  \mathcal{G}_q(S_1)
  \neq
  \mathcal{G}_q(S_2),
\end{equation}

and hence

\begin{equation}
  S_1
  \not\sim_q
  S_2.
\end{equation}

Thus,

\begin{equation}
  \operatorname{Behavior}_T(S_1)
  =
  \operatorname{Behavior}_T(S_2)
  \not\Rightarrow
  S_1\sim_q S_2.
\end{equation}

Behavioral equivalence therefore does not imply assurability equivalence.
\hfill\(\square\)

The result does not require either system to be
safer, more correct, or more trustworthy.
It establishes only that behavioral equivalence does
not determine equivalence in the structural capacity to support assurance.

\subsection{Separating Constructions}
\label{aa.subsec.formal-separations}

We now formalize the dimension-contribution criterion introduced in Section~\ref{aa.subsec.dimension-contribution}.

Let \(d\) be a candidate structural dimension.

A \emph{separating construction} for \(d\) consists of systems

\begin{equation}
  S_d^{-},
  \qquad
  S_d^{+},
\end{equation}

and comparative assurance context \(q_d\), such that the systems differ materially in \(d\), are otherwise controlled with respect to properties material to the construction, and satisfy

\begin{equation}
  \mathcal{G}_{q_d}(S_d^{-})
  \neq
  \mathcal{G}_{q_d}(S_d^{+}).
  \label{eq:general-separation}
\end{equation}

Equivalently,

\begin{equation}
  S_d^{-}
  \not\sim_{q_d}
  S_d^{+}.
\end{equation}

This non-equivalence is sufficient to establish that \(d\) contributes to assurability.

A stronger form occurs when the construction is a \emph{conservative extension}: every capability available from \(S_d^{-}\) is preserved in \(S_d^{+}\), and at least one additional claim-material capability is introduced.

Then

\begin{equation}
  \mathcal{G}_{q_d}(S_d^{-})
  \subset
  \mathcal{G}_{q_d}(S_d^{+}),
  \label{eq:directional-separation}
\end{equation}

and therefore

\begin{equation}
  S_d^{+}
  \succ_{q_d}
  S_d^{-}.
\end{equation}

The distinction is important:

\begin{equation}
  \boxed{
    \begin{array}{c}
      \text{non-equivalence establishes contribution;}
      \\
      \text{proper inclusion establishes directional dominance.}
    \end{array}
  }
  \label{eq:separation-strength}
\end{equation}

The general theory requires only the weaker result.

\subsection{Observability}
\label{aa.subsec.observability-contribution}

\begin{proposition}[Observability Contribution]
  \label{se320.prop.ObservabilityContribution}
  Observability is an assurance-relevant structural dimension.
\end{proposition}

\textit{Construction.}
Let \(S_O^{-}\) and \(S_O^{+}\) execute the same computation

\begin{equation}
  x
  \rightarrow
  h_1
  \rightarrow
  h_2
  \rightarrow
  a
\end{equation}

and produce identical external action \(a\).

In \(S_O^{-}\), only \(x\) and \(a\) are observable.

In \(S_O^{+}\), state \(h_1\) is additionally exposed through an observation mechanism that does not otherwise alter the system properties material to the construction.

Let \(C_O\) be a claim for which determining whether \(h_1\) possessed a specified property is material.

Under suitable \(K\) and \(B\), \(S_O^{+}\) enables the capability

\begin{equation}
  g_O
  =
  \text{``observe the claim-relevant property of }h_1\text{''},
\end{equation}

while \(S_O^{-}\) does not.

Therefore,

\begin{equation}
  \mathcal{G}_{q_O}(S_O^{-})
  \neq
  \mathcal{G}_{q_O}(S_O^{+}).
\end{equation}

Hence observability admits a separating construction and is assurance-relevant.
\hfill\(\square\)

If the added observation preserves every evidentiary capability of \(S_O^{-}\), then the construction is additionally directional:

\begin{equation}
  S_O^{+}
  \succ_{q_O}
  S_O^{-}.
\end{equation}

This stronger conclusion belongs to the construction, not to a universal claim that greater observability always yields greater assurability.

\subsection{Controllability}
\label{aa.subsec.controllability-contribution}

\begin{proposition}[Controllability Contribution]
  \label{se320.prop.ControllabilityContribution}
  Controllability is an assurance-relevant structural dimension.
\end{proposition}

\textit{Construction.}
Let \(S_C^{-}\) and \(S_C^{+}\) expose equivalent observations during ordinary execution.

Suppose evaluation of claim \(C_C\) requires examining system behavior under condition \(x^\ast\).

The evaluator cannot establish \(x^\ast\) in \(S_C^{-}\).

In \(S_C^{+}\), an assurance interface permits the evaluator to establish \(x^\ast\), replay the relevant state, substitute the required environmental condition, or otherwise perform the intervention without altering the property material to \(C_C\).

Thus \(S_C^{+}\) enables

\begin{equation}
  g_C
  =
  \text{``obtain evidence under condition }x^\ast\text{''},
\end{equation}

while \(S_C^{-}\) does not.

Therefore,

\begin{equation}
  \mathcal{G}_{q_C}(S_C^{-})
  \neq
  \mathcal{G}_{q_C}(S_C^{+}).
\end{equation}

Controllability is therefore assurance-relevant.
\hfill\(\square\)

When the intervention interface conservatively adds this capability while preserving the capabilities already available from \(S_C^{-}\),

\begin{equation}
  S_C^{+}
  \succ_{q_C}
  S_C^{-}.
\end{equation}

\subsection{Traceability}
\label{aa.subsec.traceability-contribution}

\begin{proposition}[Traceability Contribution]
  \label{se320.prop.TraceabilityContribution}
  Traceability is an assurance-relevant structural dimension
  distinct from mere event observability.
\end{proposition}

\textit{Construction.}
Let \(S_T^{-}\) and \(S_T^{+}\) expose the same consequential event \(a\).

In \(S_T^{-}\), the evidence establishes that \(a\) occurred but does not reliably associate it with its originating component or causal history.

In \(S_T^{+}\), authenticated provenance associates

\begin{equation}
  \text{initiator}
  \rightarrow
  \text{decision}
  \rightarrow
  \text{authorization}
  \rightarrow
  \text{tool invocation}
  \rightarrow
  a.
\end{equation}

Let \(C_T\) concern whether \(a\) originated from a specified component or authorized pathway.

Both systems expose occurrence of \(a\).
The separating capability is therefore not event observation.

System \(S_T^{+}\) additionally enables

\begin{equation}
  g_T
  =
  \text{``attribute }a\text{ to its relevant origin or pathway''}.
\end{equation}

Thus,

\begin{equation}
  \mathcal{G}_{q_T}(S_T^{-})
  \neq
  \mathcal{G}_{q_T}(S_T^{+}).
\end{equation}

Traceability is therefore assurance-relevant even when event observability is held fixed.
\hfill\(\square\)

For a conservative extension,

\begin{equation}
  S_T^{+}
  \succ_{q_T}
  S_T^{-}.
\end{equation}

\subsection{Reconstructability}
\label{aa.subsec.reconstructability-contribution}

\begin{proposition}[Reconstructability Contribution]
  \label{se320.prop.ReconstructabilityContribution}
  Reconstructability is an assurance-relevant structural dimension.
\end{proposition}

\textit{Construction.}
Let \(S_R^{-}\) and \(S_R^{+}\) have equivalent runtime behavior and equivalent live observability.

Following execution, \(S_R^{-}\) discards transient state and contextual information required to reconstruct the sequence leading to consequential action \(a\).

System \(S_R^{+}\) preserves sufficient inputs, configuration, relevant state transitions, external responses, and action history to permit such reconstruction.

Let \(C_R\) be a retrospective claim concerning the pathway by which \(a\) occurred.

During live execution, the systems may expose equivalent information.
After the event, however, \(S_R^{+}\) enables

\begin{equation}
  g_R
  =
  \text{``reconstruct the claim-relevant execution leading to }a\text{''},
\end{equation}

while \(S_R^{-}\) does not.

Therefore,

\begin{equation}
  \mathcal{G}_{q_R}(S_R^{-})
  \neq
  \mathcal{G}_{q_R}(S_R^{+}).
\end{equation}

Reconstructability is therefore assurance-relevant even when live observability is held fixed.
\hfill\(\square\)

If reconstruction evidence is added without removing other capabilities, then

\begin{equation}
  S_R^{+}
  \succ_{q_R}
  S_R^{-}.
\end{equation}

\subsection{Independence}
\label{aa.subsec.independence-contribution}

\begin{proposition}[Independence Contribution]
  \label{se320.prop.IndependenceContribution}
  Independent evidence generation or corroboration is an
  assurance-relevant structural dimension.
\end{proposition}

\textit{Construction.}
Let \(S_N^{-}\) and \(S_N^{+}\) provide the same system-generated record concerning action \(a\).

In \(S_N^{-}\), that record is the only evidence source.

In \(S_N^{+}\), an independent component outside the control of the component under evaluation produces corroborating evidence concerning the same event.

Let \(C_N\) be a claim for which failure, deception, compromise, or error in the evaluated component is within the threat model specified by \(K\).

The system-generated record may possess the same nominal integrity in both systems.

However, \(S_N^{+}\) additionally enables

\begin{equation}
  g_N
  =
  \text{``corroborate the claim-relevant event independently''}.
\end{equation}

Therefore,

\begin{equation}
  \mathcal{G}_{q_N}(S_N^{-})
  \neq
  \mathcal{G}_{q_N}(S_N^{+}).
\end{equation}

Independence is therefore assurance-relevant even when the nominal integrity of the original evidence artifact is held fixed.
\hfill\(\square\)

Where independent corroboration preserves all previously available capabilities,

\begin{equation}
  S_N^{+}
  \succ_{q_N}
  S_N^{-}.
\end{equation}

\subsection{Integrity}
\label{aa.subsec.integrity-contribution}

\begin{proposition}[Evidence Integrity Contribution]
  \label{se320.prop.EvidenceIntegrityContribution}
  Evidence integrity is an assurance-relevant structural dimension
  distinct from nominal evidence availability.
\end{proposition}

\textit{Construction.}
Let \(S_I^{-}\) and \(S_I^{+}\) generate and expose the same class of execution records.

In \(S_I^{-}\), an actor included within the threat model can silently alter, fabricate, delete, substitute, or suppress those records.

In \(S_I^{+}\), the evidence-producing environment prevents such modification or makes it detectably evident.

Let \(C_I\) concern whether an unauthorized action occurred during the recorded interval.

Both systems nominally expose the same record type.
The difference concerns whether manipulation of that historical evidence can itself be detected.

System \(S_I^{+}\) therefore enables

\begin{equation}
  g_I
  =
  \text{``detect claim-material tampering or suppression of the record''},
\end{equation}

while \(S_I^{-}\) does not.

Hence,

\begin{equation}
  \mathcal{G}_{q_I}(S_I^{-})
  \neq
  \mathcal{G}_{q_I}(S_I^{+}).
\end{equation}

Evidence integrity is therefore assurance-relevant even when nominal evidence availability is held fixed.
\hfill\(\square\)

If integrity protection adds this capability without removing another,

\begin{equation}
  S_I^{+}
  \succ_{q_I}
  S_I^{-}.
\end{equation}

\subsection{Evidence Coverage}
\label{aa.subsec.evidence-coverage-contribution}

\begin{proposition}[Evidence Coverage Contribution]
  \label{se320.prop.EvidenceCoverageContribution}
  Evidence coverage is an assurance-relevant structural dimension.
\end{proposition}

\textit{Construction.}
Let the claim-relevant action surface be

\begin{equation}
  A_C
  =
  \{
  a_1,a_2,\ldots,a_n
  \}.
\end{equation}

Let \(S_E^{-}\) and \(S_E^{+}\) provide equivalent evidence for every event observed by both systems.

System \(S_E^{-}\) exposes evidence only for

\begin{equation}
  A^{-}
  \subset
  A_C,
\end{equation}

while \(S_E^{+}\) exposes corresponding evidence across all of \(A_C\).

Choose

\begin{equation}
  a_j
  \in
  A_C
  \setminus
  A^{-},
\end{equation}

and let \(C_E\) concern whether an unauthorized action occurred anywhere in \(A_C\).

Even if observations over \(A^{-}\) are equally observable, traceable, independently corroborated, and integrity-protected in both systems, \(S_E^{-}\) provides no corresponding evidence capability over \(a_j\).

System \(S_E^{+}\) therefore enables at least one claim-material capability absent from \(S_E^{-}\):

\begin{equation}
  g_E
  =
  \text{``investigate occurrence of unauthorized action }a_j\text{''}.
\end{equation}

Thus,

\begin{equation}
  \mathcal{G}_{q_E}(S_E^{-})
  \neq
  \mathcal{G}_{q_E}(S_E^{+}).
\end{equation}

Evidence coverage is therefore assurance-relevant independently of the quality of evidence over the already-covered surface.
\hfill\(\square\)

If coverage is expanded without losing existing capabilities,

\begin{equation}
  S_E^{+}
  \succ_{q_E}
  S_E^{-}.
\end{equation}

\subsection{What the Separating Constructions Establish}
\label{aa.subsec.what-separation-establishes}

Propositions~2--8 present candidate separating arguments for the
seven structural dimensions.

A completed separating construction has two obligations:
a difference in the designated structural profile, and a
claim-material difference in the evidentiary capabilities
available under the specified assurance context.

The second obligation is not a consequence of the first.
It must be established separately.

The companion Lean~4 formalization, version 0.4.1,
provides dimension-specific, profile-grounded constructions for
observability, evidence coverage, traceability, and reconstructability.
These constructions require capability non-equivalence as a
separate hypothesis.

Controllability, independence, and integrity remain supported
by the informal arguments presented here.

For any candidate dimension \(d\), a completed separation requires

\begin{equation}
  \mathcal{G}_{q_d}(S_d^{-})
  \neq
  \mathcal{G}_{q_d}(S_d^{+}).
  \label{eq:all-dimensions-separate}
\end{equation}

Such a result establishes contribution in the specified context.
It does not establish universal monotonicity:

\begin{equation}
  \forall q,\;
  d(S_1)>d(S_2)
  \Rightarrow
  S_1\succ_q S_2.
  \label{eq:no-universal-monotonicity}
\end{equation}

Indeed, the notation \(d(S_1)>d(S_2)\) may itself be
inappropriate for dimensions without a natural numerical scale.

Nor does a separating construction establish that the dimension
is required for every assurance claim.

The contribution criterion is existential: a dimension matters
to comparative assurability when a controlled structural difference
is accompanied by an independently justified difference in
claim-material evidentiary capabilities.

The formalization therefore provides a basis for testing
candidate dimensions without asserting completeness,
independence, or universal necessity.

\subsection{Contribution Is Not Sufficiency}
\label{aa.subsec.formal-non-sufficiency}

The preceding propositions should not be combined into a sufficiency claim.

Let

\begin{equation}
  D_S^{(0)}
  =
  \{
  O,C,T,R,N,I,E
  \}
\end{equation}

denote the current candidate structural set.

For each candidate dimension, establishing contribution requires
a separating construction satisfying

\begin{equation}
  \exists q_d,S_d^{-},S_d^{+}
  :
  S_d^{-}
  \not\sim_{q_d}
  S_d^{+}.
\end{equation}

This obligation includes an independently justified capability
difference.

Even when satisfied, these obligations do not establish

\begin{equation}
  D_S^{(0)}
  \Longrightarrow
  \text{complete characterization of assurability},
\end{equation}

and they do not establish

\begin{equation}
  \bigwedge_{d\in D_S^{(0)}} d
  \Longrightarrow
  \text{sufficient assurance}.
  \label{eq:no-joint-sufficiency}
\end{equation}

The first unsupported implication would assert completeness of the structural model.

The second would conflate structural evidence capacity with the downstream truth or justification of an assurance claim.

Neither assertion is required by the theory.

\subsection{Context Dependence}
\label{aa.subsec.context-dependence}

The same structural difference may affect assurability under one comparative assurance context while being irrelevant under another.

For example, additional internal observability may introduce
important evidentiary capabilities for a claim concerning internal computational state
while providing no additional capability for a claim resolvable entirely through independently observable external effects.

Likewise, independent corroboration may substantially affect assurability when the original evidence source
falls within the threat model, while contributing no material distinction under a context that assumes the source to be trusted.

Thus,

\begin{equation}
  S_1\succ_{q_1}S_2
\end{equation}

does not imply

\begin{equation}
  S_1\succ_{q_2}S_2.
\end{equation}

Nor does

\begin{equation}
  S_1\not\sim_{q_1}S_2
\end{equation}

imply

\begin{equation}
  S_1\not\sim_{q_2}S_2.
\end{equation}

Assurability comparison is inherently relative to the assurance claim and context.

\subsection{Access Expansion}
\label{aa.subsec.access-monotonicity}

Access conditions affect obtainable evidence.

Let

\begin{equation}
  K_1
  \preceq_K
  K_2
\end{equation}

mean that every evidence source accessible under \(K_1\) remains accessible under \(K_2\), while \(K_2\) may provide additional access.

Holding \(C\), \(B\), the system, and applicable inferential standards fixed,
suppose the additional access can be exercised within \(B\).

Then

\begin{equation}
  \mathcal{E}_{(C,K_1,B)}(S)
  \subseteq
  \mathcal{E}_{(C,K_2,B)}(S).
  \label{eq:access-evidence-monotonicity}
\end{equation}

Access expansion therefore cannot reduce the set of obtainable evidence under these assumptions.

It does not follow that

\begin{equation}
  \mathcal{G}_{(C,K_1,B)}(S)
  \subset
  \mathcal{G}_{(C,K_2,B)}(S).
\end{equation}

Additional evidence may enable no additional claim-material capability.
It may be irrelevant, redundant, invalid for the intended inference,
dependent upon an existing source, or otherwise incapable of distinguishing alternatives material to \(C\).

Therefore,

\begin{equation}
  \boxed{
    \text{more access}
    \Rightarrow
    \text{no less obtainable evidence}
  }
\end{equation}

under the stated assumptions, but

\begin{equation}
  \boxed{
    \text{more obtainable evidence}
    \not\Rightarrow
    \text{strictly greater assurability}.
  }
\end{equation}

This distinction prevents telemetry volume from being confused with structural assurance capability.

\subsection{Resource Dependence}
\label{aa.subsec.resource-dependence}

Structural capabilities also interact with resource constraints.

Let

\begin{equation}
  \operatorname{cost}(e,S,K)
\end{equation}

represent the resources required to obtain evidence \(e\) from or about \(S\) under access conditions \(K\).

The quantity need not be monetary.
It may represent computation, time, evaluator effort, storage, latency, specialized expertise,
disclosure burden, operational disruption, or another resource constrained by \(B\).

Evidence \(e\) belongs to the practically obtainable set only when its acquisition satisfies the applicable resource constraint:

\begin{equation}
  \operatorname{cost}(e,S,K)
  \preceq
  B.
\end{equation}

Consequently, two systems may make the same type of evidence theoretically accessible while differing in practical assurability.

Suppose evidence capability \(g\) requires evidence \(e\), and

\begin{equation}
  \operatorname{cost}(e,S_1,K)
  \preceq
  B
\end{equation}

while

\begin{equation}
  \operatorname{cost}(e,S_2,K)
  \npreceq
  B.
\end{equation}

Then \(g\) may belong to

\begin{equation}
  \mathcal{G}_q(S_1)
\end{equation}

while being absent from

\begin{equation}
  \mathcal{G}_q(S_2).
\end{equation}

Thus, architectural differences that make evidence inexpensive, routine, and
operationally useful can affect assurability even when comparable evidence
is theoretically recoverable from another system given unconstrained resources.

Resource constraints are therefore part of the assurance context rather than an incidental implementation concern.

\subsection{Dominance and Incomparability}
\label{aa.subsec.dominance}

The capability-set formulation makes dominance explicit.

If

\begin{equation}
  \mathcal{G}_q(S_2)
  \subset
  \mathcal{G}_q(S_1),
\end{equation}

then

\begin{equation}
  S_1
  \succ_q
  S_2.
\end{equation}

System \(S_1\) is a strict assurability extension of \(S_2\) for that context:
it preserves every claim-material evidentiary capability available from \(S_2\) and adds at least one more.

However, suppose

\begin{equation}
  g_a
  \in
  \mathcal{G}_q(S_1)
  \setminus
  \mathcal{G}_q(S_2)
\end{equation}

and

\begin{equation}
  g_b
  \in
  \mathcal{G}_q(S_2)
  \setminus
  \mathcal{G}_q(S_1).
\end{equation}

Then neither capability set contains the other, and

\begin{equation}
  S_1
  \parallel_q
  S_2.
\end{equation}

For example, one architecture may enable detailed internal-state examination
while another provides stronger independently generated evidence of externally consequential actions.

If both capabilities are material to \(C\),
the general framework need not decide that one architecture is globally superior.

The systems are incomparable under \(q\).

The result indicates an assurance tradeoff.

\subsection{Why No Universal Scalar Score Is Required}
\label{aa.subsec.no-scalar}

The framework therefore does not define a quantity such as

\begin{equation}
  \mathcal{A}(S)=0.87.
\end{equation}

Such a scalar would require additional commitments that have not been established:

\begin{enumerate}
  \item a complete set of assurance-relevant structural dimensions;
  \item valid measurement scales for those dimensions;
  \item a mapping from dimensions to evidentiary capabilities;
  \item a weighting or aggregation rule for heterogeneous capabilities; and
  \item a justification for the tradeoffs implied by that aggregation.
\end{enumerate}

None is required for comparative assurability.

Set containment already permits rigorous statements of equivalence and dominance:

\begin{equation}
  \mathcal{G}_q(S_1)
  =
  \mathcal{G}_q(S_2),
\end{equation}

\begin{equation}
  \mathcal{G}_q(S_1)
  \subset
  \mathcal{G}_q(S_2),
\end{equation}

or incomparability when neither inclusion holds.

A scalar aggregation could later be introduced for a specific application, assurance authority, procurement regime, or engineering objective.

Such an aggregation would be an additional decision model, not part of the foundational definition of assurability.

\subsection{Derived Capabilities and Primitive Dimensions}
\label{aa.subsec.derived-formal}

The distinction between structural dimensions and evidentiary capabilities
also clarifies why properties such as reproducibility, auditability, monitorability, testability, and diagnosability
need not initially be treated as primitive dimensions.

A derived capability may depend upon several structural dimensions.

For example, a capability

\begin{equation}
  g_{\mathrm{replay}}
  =
  \text{``reproduce a relevant execution''}
\end{equation}

may require some combination of:

\begin{equation}
  \text{traceability},
  \quad
  \text{reconstructability},
  \quad
  \text{controllability},
  \quad
  \text{integrity},
\end{equation}

together with preservation of configuration and environmental state.

The formal model therefore need not assign one primitive dimension to every familiar assurance concept.

Structural dimensions characterize properties of the evidence-producing architecture.

Evidentiary capabilities characterize what that structure makes possible for a particular assurance problem.

The mapping

\begin{equation}
  D_S
  \longrightarrow
  \mathcal{E}_q(S)
  \longrightarrow
  \mathcal{G}_q(S)
  \label{eq:dimensions-to-capabilities}
\end{equation}

may be many-to-many.

This distinction also supplies a test for whether a proposed new concept
should be treated as a primitive structural dimension or as a capability derived from existing dimensions.

\subsection{Extension of the Structural Model}
\label{aa.subsec.formal-open-endedness}

Let the current structural candidate set be

\begin{equation}
  D_S^{(0)}
  =
  \{
  O,C,T,R,N,I,E
  \}.
\end{equation}

Suppose future work identifies structural property \(d^\ast\) not represented by \(D_S^{(0)}\).

The property should be admitted as an additional assurance-relevant dimension if there exists a separating construction

\begin{equation}
  S_{d^\ast}^{-},
  \qquad
  S_{d^\ast}^{+},
  \qquad
  q_{d^\ast}
\end{equation}

such that

\begin{equation}
  S_{d^\ast}^{-}
  \not\sim_{q_{d^\ast}}
  S_{d^\ast}^{+},
\end{equation}

and the resulting capability difference cannot already be represented by variation in the existing dimensions.

The structural set can then be extended:

\begin{equation}
  D_S^{(1)}
  =
  D_S^{(0)}
  \cup
  \{d^\ast\}.
\end{equation}

Conversely, if a proposed dimension produces no separating construction independent of the existing dimensions,
its treatment as a distinct primitive dimension should be reconsidered.

This yields two complementary tests:

\begin{quote}
  \textbf{Admission test.}
  Does the proposed dimension expose an assurance-relevant capability distinction
  that the current structural model cannot represent?
\end{quote}

\begin{quote}
  \textbf{Reduction test.}
  Can the apparent contribution of the proposed dimension be represented completely
  through existing structural dimensions?
\end{quote}

The first supports extension.

The second supports simplification.

The theory therefore remains explicitly revisable as assurance science and AI architectures evolve.

\subsection{Structural Neutrality}
\label{aa.subsec.formal-neutrality}

The capability ordering does not encode whether an assurance claim receives a favorable outcome.

Suppose

\begin{equation}
  S_1
  \succ_q
  S_2.
\end{equation}

This means that \(S_1\) enables a strict superset of the claim-material evidentiary capabilities available from \(S_2\).

It does not imply

\begin{equation}
  C(S_1)
\end{equation}

is more likely to be true than

\begin{equation}
  C(S_2).
\end{equation}

Indeed, the additional capability may expose evidence that defeats \(C\).

Thus,

\begin{equation}
  \boxed{
    S_1\succ_q S_2
    \not\Rightarrow
    \text{a more favorable assurance conclusion for }S_1.
  }
\end{equation}

This is an essential neutrality property of the framework.

Structural assurability describes the capacity for legitimate examination.
It does not encode the desired result of that examination.

\subsection{Summary of Formal Claims}
\label{aa.subsec.formal-summary}

The formal contribution of this section is intentionally bounded.

We establish the following.

\begin{enumerate}

  \item For a fixed comparative assurance
        context \(q=(C,K,B)\),
        a system has a set of practically obtainable evidence
        \[
          \mathcal{E}_q(S),
        \]
        a claim-material subset
        \[
          \mathcal{E}^{C}_q(S),
        \]
        and a corresponding set of evidentiary capabilities
        \[
          \mathcal{G}_q(S).
        \]

  \item Evidence sets are compared through the claim-relative capability map
        \[
          \Gamma_C(E),
        \]
        rather than through raw evidence volume or artifact identity.

  \item Comparative assurability is defined by capability containment:
        \[
          S_1\succeq_q S_2
          \quad\Longleftrightarrow\quad
          \mathcal{G}_q(S_2)
          \subseteq
          \mathcal{G}_q(S_1).
        \]

  \item Behavioral equivalence does not imply assurability equivalence.

  \item A candidate structural dimension contributes to assurability when it admits a separating construction for which
        \[
          S_d^{-}
          \not\sim_{q_d}
          S_d^{+}.
        \]

  \item Strict dominance,
        \[
          S_d^{+}\succ_{q_d}S_d^{-},
        \]
        is a stronger construction-specific result requiring preservation
        of the weaker system's claim-material capabilities together with at least one additional capability.

  \item Seven candidate dimensions are examined through illustrative
        separating arguments. Four have corresponding
        dimension-specific, profile-grounded constructions in
        the companion Lean~4 formalization.

  \item Establishing contribution requires both structural-profile
        separation and capability non-equivalence.
        Neither obligation establishes universal necessity,
        completeness, or sufficiency.

  \item The assurability relation permits equivalence, strict dominance, and incomparability without requiring a universal scalar score.

  \item Additional structural dimensions may be admitted when they expose assurance-relevant capability distinctions not representable by the existing dimensions.

\end{enumerate}

The resulting theory separates three questions that are frequently collapsed:

\begin{equation}
  \boxed{
    \begin{array}{c}
      \text{What evidence does the structure make obtainable?}
      \\[4pt]
      \downarrow
      \\[4pt]
      \text{What assurance-relevant capabilities does that evidence enable?}
      \\[4pt]
      \downarrow
      \\[4pt]
      \text{What conclusion should an assurance argument draw?}
    \end{array}
  }
\end{equation}

Structural assurability addresses the first two while deliberately leaving the final inferential judgment to the assurance process.

The following section applies the framework to representative AI architectures to determine
whether systems with similar functional capabilities can occupy different positions in the comparative assurability relation.

% !TeX root = ../se320_structural_assurability.tex

% ==================================================
\section{Applying Structural Assurability to AI Architectures}
\label{aa.sec.applications}
% ==================================================

The preceding sections developed structural assurability as a claim-relative
property concerning the capacity of a system and its evidence-producing
environment to make assurance-relevant distinctions possible.
This section applies that framework to representative AI architectures.

The purpose is not to evaluate particular commercial systems or prescribe
a preferred architecture.
Instead, the examples examine whether the formal distinctions developed in
Sections~\ref{aa.sec.foundations} and~\ref{aa.sec.formal-model}
reveal meaningful differences among systems that may otherwise appear
functionally similar.

Three cases are considered.

The first demonstrates strict comparative dominance for a claim concerning
unauthorized consequential action.

The second demonstrates assurability equivalence for a claim that can be
resolved primarily through externally observable behavior.

The third demonstrates incomparability when different architectures enable
different assurance-relevant distinctions.

Together, the cases illustrate that structural assurability is claim-relative
and need not induce a universal ranking of architectures.

\subsection{Representative Architectures}
\label{aa.subsec.representative-architectures}

Consider three abstract AI system architectures.

They are intentionally simplified.
Each represents a class of evidence-producing structure rather than a
specific implementation.

\subsubsection{Opaque Service Architecture}

Let \(S_O\) denote an AI system exposed primarily through an input-output
interface.

The system accepts requests and returns outputs.
Standard application-level records may establish that a request occurred
and that a response was returned, but the assurance interface exposes little
or no trusted evidence concerning internal execution, tool use,
authorization decisions, retrieved information, intermediate state,
or causal provenance.

A simplified evidence surface is

\begin{equation}
  S_O:
  \qquad
  \{
  \text{input},
  \text{output},
  \text{request metadata}
  \}.
\end{equation}

This architecture is not assumed to be unsafe or unreliable.
It is simply structurally limited in the assurance-relevant distinctions
that can be made from its exposed evidence.

\subsubsection{Instrumented Agent Architecture}

Let \(S_A\) denote an agentic AI system with runtime instrumentation.

In addition to externally visible inputs and outputs, the system records
evidence concerning consequential execution events, such as

\begin{equation}
  \begin{aligned}
    S_A:\quad
    \{ &
    \text{input},
    \text{output},
    \text{tool invocations},
    \text{retrieval events}, \\
       &
    \text{authorization decisions},
    \text{memory updates},
    \text{state transitions},
    \text{provenance}
    \}.
  \end{aligned}
\end{equation}

The architecture therefore provides a larger evidence surface concerning
system execution.

However, the evidence may still be generated and preserved within
infrastructure controlled by the same operator or execution environment
as the agent itself.

\subsubsection{Assurability-Oriented Architecture}

Let \(S_R\) denote an architecture explicitly designed to preserve
assurance-relevant evidence.

It includes the runtime evidence available from \(S_A\) and adds
mechanisms such as

\begin{equation}
  \begin{aligned}
    S_R:\quad \{ &
    \text{instrumented execution evidence},
    \text{independent action observation},   \\
                 &
    \text{tamper-evident evidence preservation},
    \text{authenticated provenance},         \\
                 &
    \text{controlled replay or intervention},
    \text{explicit action-surface coverage}, \\
                 &
    \text{selected internal-state instrumentation}
    \}.
  \end{aligned}
\end{equation}

The term \emph{assurability-oriented} is descriptive rather than normative.
The architecture is designed to increase the range of assurance-relevant
distinctions that can be made under specified conditions.

It does not follow that the system satisfies any particular assurance claim.

\subsection{Case 1: Unauthorized Consequential Action}
\label{aa.subsec.case-unauthorized-action}

Consider the assurance claim

\begin{equation}
  \begin{aligned}
    C_1 ={} & \text{``The agent performed no unauthorized consequential action} \\
            & \text{during interval }T\text{.''}
  \end{aligned}
  \label{eq:claim-unauthorized-action}
\end{equation}

Suppose a consequential external action \(a\) is observed during \(T\).

Relevant assurance questions include:

\begin{align*}
  Q_1 & = \text{Did action }a\text{ occur?}                                                  \\
  Q_2 & = \text{Which component or actor initiated }a\text{?}                                \\
  Q_3 & = \text{What authorization state applied when }a\text{ was initiated?}               \\
  Q_4 & = \text{Through what execution pathway did }a\text{ occur?}                          \\
  Q_5 & = \text{Could the evidence concerning }a\text{ have been altered or suppressed?}     \\
  Q_6 & = \text{Is the evidence concerning }a\text{ corroborated by a source}                \\
      & \quad\text{outside the evaluated component?}                                         \\                                                        \\
  Q_7 & = \text{Were all consequential-action channels represented by the evidence surface?}
\end{align*}

These questions define claim-relevant distinctions.

For example, \(Q_2\) may require distinguishing

\begin{equation}
  \text{agent initiated }a
  \quad\text{from}\quad
  \text{human initiated }a,
\end{equation}

while \(Q_3\) may require distinguishing

\begin{equation}
  \text{authorized action}
  \quad\text{from}\quad
  \text{unauthorized action}.
\end{equation}

\subsubsection{Opaque Service}

Under \(S_O\), external evidence may establish that a result occurred,
but the evidence surface may not permit reliable attribution of the
action to the agent, a human operator, an external service, or another
component.

Thus, the architecture may permit resolution of \(Q_1\) while leaving
\(Q_2\)--\(Q_7\) substantially unresolved.

Its claim-relevant distinction set may therefore be represented
schematically as

\begin{equation}
  \mathcal{G}_{q_1}(S_O)
  =
  \{
  g_{\mathrm{occurrence}}
  \}.
\end{equation}

\subsubsection{Instrumented Agent}

System \(S_A\) additionally exposes runtime traces linking actions to
tool invocations, authorization decisions, retrieval events, and state
transitions.

It may therefore permit distinctions concerning:

\begin{equation}
  \begin{aligned}
     & \text{action occurrence},   \\
     & \text{action initiator},    \\
     & \text{authorization state}, \\
     & \text{execution pathway}.
  \end{aligned}
\end{equation}

Schematically,

\begin{equation}
  \mathcal{G}_{q_1}(S_A)
  =
  \{
  g_{\mathrm{occurrence}},
  g_{\mathrm{origin}},
  g_{\mathrm{authorization}},
  g_{\mathrm{pathway}}
  \}.
\end{equation}

However, if all records are generated and preserved within the same
administrative or computational boundary as the agent, questions
concerning evidence manipulation or independent corroboration may
remain unresolved.

\subsubsection{Assurability-Oriented Architecture}

System \(S_R\) preserves the distinctions available from \(S_A\) while
adding independently generated observations, tamper-evident records,
authenticated provenance, and explicit monitoring of consequential-action
channels.

It may therefore additionally distinguish:

\begin{equation}
  \text{unaltered evidence}
  \quad\text{from}\quad
  \text{evidence subject to undetected manipulation},
\end{equation}

and

\begin{equation}
  \text{complete monitored action surface}
  \quad\text{from}\quad
  \text{partially observed action surface}.
\end{equation}

Schematically,

\begin{equation}
  \begin{aligned}
    \mathcal{G}_{q_1}(S_R)
    =
    \{ &
    g_{\mathrm{occurrence}},
    g_{\mathrm{origin}},
    g_{\mathrm{authorization}},
    g_{\mathrm{pathway}}, \\
       &
    g_{\mathrm{integrity}},
    g_{\mathrm{coverage}},
    g_{\mathrm{independent\ corroboration}}
    \}.
  \end{aligned}
\end{equation}

If each architecture preserves the capabilities of the preceding
architecture while adding the capabilities shown above, then

\begin{equation}
  \mathcal{G}_{q_1}(S_O)
  \subset
  \mathcal{G}_{q_1}(S_A)
  \subset
  \mathcal{G}_{q_1}(S_R).
\end{equation}

Therefore,

\begin{equation}
  S_R
  \succ_{q_1}
  S_A
  \succ_{q_1}
  S_O.
  \label{eq:case1-order}
\end{equation}

This ordering is not a general ranking of the three architectures.

It holds only for the specified claim and assurance context.

The result illustrates how structural differences can create strict
comparative assurability without establishing that any of the systems
actually satisfies \(C_1\).

Indeed, \(S_R\) may provide stronger evidence that an unauthorized action
\emph{did} occur.

\subsection{Case 2: Externally Observable Output Property}
\label{aa.subsec.case-output-property}

Consider a different claim:

\begin{equation}
  \begin{aligned}
    C_2 ={} & \text{``The response returned for request }x\text{ contains no} \\
            & \text{instance of prohibited output class }P\text{.''}
  \end{aligned}
  \label{eq:claim-output}
\end{equation}

Assume that membership in \(P\) can be determined completely and
reliably from the externally returned response.

The material assurance question is then simply whether

\begin{equation}
  y \in P,
\end{equation}

where \(y\) is the externally observable output.

All three architectures expose \(y\).

Additional information concerning tool calls, internal execution,
provenance, independent runtime observation, or state transitions may
be valuable for other purposes, but under the assumptions of this claim
it enables no additional distinction required to determine whether
\(y\in P\).

Thus,

\begin{equation}
  \mathcal{G}_{q_2}(S_O)
  =
  \mathcal{G}_{q_2}(S_A)
  =
  \mathcal{G}_{q_2}(S_R).
\end{equation}

Therefore,

\begin{equation}
  S_O
  \sim_{q_2}
  S_A
  \sim_{q_2}
  S_R.
  \label{eq:case2-equivalence}
\end{equation}

This case is important because it prevents a common misinterpretation
of structural assurability.

More instrumentation does not necessarily imply greater assurability.

If the assurance-relevant distinction can already be resolved completely
from externally available evidence, additional structural evidence may
contribute nothing to that particular assurance context.

Thus,

\begin{equation}
  S_R\succ_{q_1}S_O
\end{equation}

and

\begin{equation}
  S_R\sim_{q_2}S_O
\end{equation}

can both hold without contradiction.

Assurability is claim-relative.

\subsection{Case 3: Internal Visibility Versus Independent Evidence}
\label{aa.subsec.case-incomparability}

The previous examples produced strict dominance and equivalence.

A third case illustrates incomparability.

Consider two architectures, \(S_V\) and \(S_N\).

System \(S_V\) provides extensive visibility into internal execution.
It exposes selected intermediate representations, internal state transitions,
detailed tool-selection processes, and model-generated explanations of
relevant actions.

However, these observations are produced and preserved by infrastructure
controlled by the same operator and execution environment as the system
under evaluation.

System \(S_N\), by contrast, exposes little internal state.
It provides only limited visibility into model computation.

However, consequential external actions are independently observed
by a separate enforcement or monitoring component.
Those records are authenticated and protected from retrospective
alteration by the evaluated system.

Consider the claim

\begin{equation}
  C_3 =
  \text{``The agent intentionally attempted an unauthorized consequential action.''}
  \label{eq:claim-intention-action}
\end{equation}

The claim contains at least two materially different components.

One concerns internal process:

\begin{equation}
  \begin{aligned}
    q_{\mathrm{internal}} ={} & \text{Did the agent enter a state or process}     \\
                              & \text{relevant to attempted unauthorized action?}
  \end{aligned}
\end{equation}

The other concerns external execution and attribution:

\begin{equation}
  \begin{aligned}
    q_{\mathrm{external}} ={} & \text{Did a consequential unauthorized action occur,} \\
                              & \text{and can it be independently attributed?}
  \end{aligned}
\end{equation}

System \(S_V\) may enable distinctions over internal computational
behavior unavailable from \(S_N\):

\begin{equation}
  g_{\mathrm{internal}}
  \in
  \mathcal{G}_{q_3}(S_V)
  \setminus
  \mathcal{G}_{q_3}(S_N).
\end{equation}

System \(S_N\), however, may enable independently corroborated
distinctions concerning consequential external action that \(S_V\)
cannot provide:

\begin{equation}
  g_{\mathrm{independent}}
  \in
  \mathcal{G}_{q_3}(S_N)
  \setminus
  \mathcal{G}_{q_3}(S_V).
\end{equation}

Accordingly,

\begin{equation}
  \mathcal{G}_{q_3}(S_V)
  \nsubseteq
  \mathcal{G}_{q_3}(S_N)
\end{equation}

and

\begin{equation}
  \mathcal{G}_{q_3}(S_N)
  \nsubseteq
  \mathcal{G}_{q_3}(S_V).
\end{equation}

Therefore,

\begin{equation}
  S_V
  \parallel_{q_3}
  S_N.
  \label{eq:case3-incomparable}
\end{equation}

Neither system dominates.

A scalar assurance score would require assigning a relative value
to internal visibility and independent external evidence.

The structural framework does not impose such a tradeoff.

That decision belongs to the specific assurance process, threat model,
or decision authority.

\subsection{Recasting the Cases as Claim-Relevant Distinctions}
\label{aa.subsec.case-distinctions}

The three cases can be interpreted through a common principle.

Let \(\Omega_C\) represent alternatives relevant to assurance claim \(C\).

Evidence contributes to structural assurability when it permits
distinctions among alternatives in \(\Omega_C\) that would otherwise
remain unresolved.

For Case~1, relevant alternatives include:

\begin{equation}
  \begin{aligned}
    \Omega_{C_1}
    =
    \{ &
    \text{no consequential action},     \\
       &
    \text{authorized agent action},     \\
       &
    \text{unauthorized agent action},   \\
       &
    \text{human-initiated action},      \\
       &
    \text{externally initiated action}, \\
       &
    \text{action with unreliable historical record}
    \}.
  \end{aligned}
\end{equation}

The richer architectures expose evidence permitting progressively more
of these alternatives to be distinguished.

For Case~2,

\begin{equation}
  \Omega_{C_2}
  =
  \{
  y\in P,
  y\notin P
  \}.
\end{equation}

The externally visible response already permits the relevant distinction.
Additional structural instrumentation need not increase assurability.

For Case~3, the relevant alternatives span both internal-process
hypotheses and externally observable action hypotheses.
Different architectures distinguish different subsets, producing
incomparability.

These examples motivate a more general interpretation:

\begin{quote}
  Structural assurability concerns the capacity of a system and
  its evidence-producing environment to make distinctions among
  alternatives material to an assurance claim.
\end{quote}

Under this interpretation, the structural dimensions introduced in
Section~\ref{aa.subsec.candidate-dimensions} are important because
removing one can make an assurance-relevant distinction unavailable.

\subsection{Relationship to the Structural Dimensions}
\label{aa.subsec.cases-dimensions}

The examples also illustrate how the candidate structural dimensions
affect assurance-relevant distinctions.

\paragraph{Observability.}
Observability determines whether relevant states, events, or effects
enter the available evidence surface at all.

\paragraph{Controllability.}
Controllability determines whether an evaluator can create or reproduce
conditions required to distinguish claim-relevant alternatives.

\paragraph{Traceability.}
Traceability permits observed events to be distinguished according
to origin, causal pathway, or transformation history.

\paragraph{Reconstructability.}
Reconstructability preserves distinctions that may be required after
the original execution has ended.

\paragraph{Independence.}
Independence permits distinctions between system self-report and
corroboration generated outside the component or actor being evaluated.

\paragraph{Integrity.}
Integrity permits distinctions between trustworthy preservation of
historical evidence and histories subject to undetected alteration
or suppression.

\paragraph{Evidence Coverage.}
Evidence coverage determines which portions of the claim-relevant state,
event, or action surface admit evidentiary distinctions at all.

The dimensions therefore need not be interpreted as components of
a numerical assurability score.

Their more fundamental role is to govern which claim-relevant
distinctions the evidence-producing architecture makes possible.

\subsection{Cost and Practical Assurability}
\label{aa.subsec.case-cost}

The examples above can also be extended to include resource constraints.

Suppose \(S_A\) and \(S_R\) both permit reconstruction of a consequential
execution.

In \(S_A\), reconstruction requires several days of specialist analysis
across heterogeneous logs.

In \(S_R\), authenticated provenance and preserved execution state permit
the same claim-relevant distinction to be established automatically
within minutes.

Under an unconstrained budget, both systems may enable the same
evidentiary capability.

Under a resource constraint \(B^\ast\),

\begin{equation}
  \operatorname{cost}
  \left(
  g_{\mathrm{reconstruction}},
  S_A,
  K
  \right)
  >
  B^\ast
\end{equation}

while

\begin{equation}
  \operatorname{cost}
  \left(
  g_{\mathrm{reconstruction}},
  S_R,
  K
  \right)
  \le
  B^\ast.
\end{equation}

Then

\begin{equation}
  g_{\mathrm{reconstruction}}
  \notin
  \mathcal{G}_{(C,K,B^\ast)}(S_A)
\end{equation}

while

\begin{equation}
  g_{\mathrm{reconstruction}}
  \in
  \mathcal{G}_{(C,K,B^\ast)}(S_R).
\end{equation}

The systems therefore differ in practical assurability despite
admitting similar evidence in principle.

This distinction is particularly important for deployed AI systems.
Assurance mechanisms that are technically possible but prohibitively
expensive, slow, fragile, or operationally disruptive may provide
little practical assurance under realistic constraints.

Structural features that make useful evidence inexpensive and routine
can therefore expand practical assurability without changing the
underlying assurance claim.

\subsection{What the Cases Demonstrate}
\label{aa.subsec.case-summary}

The three cases establish several properties of the framework.

First, systems exhibiting similar functional behavior can differ
materially in the assurance-relevant distinctions supported by their
evidence-producing structures.

Second, additional instrumentation does not automatically increase
assurability.
Its value depends upon the assurance claim.

Third, structural assurability is not a universal ranking of architectures.
The same pair of systems may exhibit strict dominance under one claim
and equivalence under another.

Fourth, architectures may be incomparable when each exposes
claim-material distinctions unavailable from the other.

Fifth, independence, integrity, provenance, internal visibility, and
coverage are not interchangeable forms of ``more evidence.''
They can enable qualitatively different assurance-relevant distinctions.

Sixth, resource constraints can convert differences in the efficiency
of evidence acquisition into differences in practical assurability.

Most importantly, the cases support the central interpretation
developed in this work:

\begin{quote}
  A system is structurally more assurable in a specified assurance
  context when its architecture and evidence-producing environment
  permit assurance-relevant distinctions that cannot practically
  be made for the comparison system, while preserving the
  distinctions already available from that system.
\end{quote}

This statement remains deliberately neutral regarding the eventual
assurance outcome.

A richer evidence structure may make it easier to support an
assurance claim, easier to refute it, or both.

Structural assurability concerns the capacity for examination,
not the conclusion of that examination.

The examples therefore reinforce the separation

\begin{equation}
  \boxed{
    \begin{aligned}
      \text{System Structure}
       & \rightarrow
      \text{Obtainable Evidence} \\
       & \rightarrow
      \text{Assurance-Relevant Distinctions}
      \rightarrow
      \text{Assurance Reasoning}.
    \end{aligned}
  }
\end{equation}

The present framework addresses the first three terms.

The final inference from evidence and distinctions to an assurance
conclusion remains the responsibility of the applicable assurance methodology.

% !TeX root = ../se320_structural_assurability.tex
% ==================================================
\section{Theoretical Results}
\label{aa.sec.theoretical-results}
% ==================================================

The definitions and results of Sections~\ref{aa.sec.foundations} and~\ref{aa.sec.formal-model}
are developed in a Lean~4 formalization.
The companion Lean~4 formalization at version 0.4.1
is authoritative for formal statements and proofs
\citep{case2026sestructuralassurability}.

\subsection{Established in the Formalization}

\begin{itemize}
  \item \emph{Resolution bound.}
        If two admissible worlds are observationally indistinguishable under an observation model
        and disagree on a claim, the model does not resolve the claim.

  \item \emph{Capability comparison.}
        Capability sets are compared by containment.
        The capability model is abstract.
        Monotonicity of capabilities under added evidence is an explicit assumption, not a derived result.

  \item \emph{Access and resource expansion.}
        The formalization states results on expanding evaluator access and resource bounds,
        and a conservative-extension result.

  \item \emph{Separating constructions.}
        A generic separating construction is stated,
        together with dimension-specific constructions for
        observability, evidence coverage, traceability, and reconstructability.
        Each construction involves two obligations:
        the two systems differ in the named dimension's profile,
        and their capability sets are not equivalent.
        The second is a separate hypothesis, not a consequence of the first.
\end{itemize}

\subsection{Not Established}

\begin{itemize}
  \item Independence, integrity, and controllability are defined as candidate dimensions
        but have no dimension-specific separating construction.
  \item The separating-construction criterion establishes contribution to assurance-relevant capability,
        not necessity or sufficiency.
  \item The capability model is not derived from claim resolution beyond the worked examples.
  \item The finite-model experiments on observability gaps, integrity, and collusion are exploratory
        and are not part of the formal theory.
\end{itemize}

% !TeX root = ../se320_structural_assurability.tex
% ==================================================
\section{Illustrative Application}
\label{aa.sec.pilot}
% ==================================================

A companion pilot repository applies the resolution bound to an external evaluation framework.
Its initial study uses the query-violation example in
NIST's initial public draft of the TEVV-Athlon framework,
NIST AI 200-2~\citep{nist2026tevvathlon}.
The pilot adds two researcher-defined claims to that example.
NIST does not assert either claim.

In a finite model where organic behavior in four scenarios varies independently,
observations from the three tested scenarios do not resolve a claim that also covers the fourth,
with or without Red Teaming observations.
If organic and adversarial behavior vary independently,
Red Teaming alone does not resolve an organic-use claim on the three tested scenarios.
Each result comes with a pair of indistinguishable worlds that disagree on the claim.

Because the model admits every combination of scenario outcomes,
these results follow from its admissibility assumptions.
They illustrate the resolution bound; they do not test the theory.
They make no claim about operational AI systems.
Comparison of evaluator contexts by capability-set containment has not yet been carried out.
The pilot record remains open.

% !TeX root = ../se320_structural_assurability.tex
% ==================================================
\section{Discussion}
\label{aa.sec.discussion}
% ==================================================

Structural assurability is claim-relative.
The same system can support one claim well and another poorly.
Comparison is by containment of capability sets, so incomparability is a possible and informative outcome.

The framework complements assurance cases, auditability, monitorability, and evaluation-rigor methods.
It addresses what evidence an architecture and its environment make obtainable.
It does not address how that evidence is argued or how much confidence it warrants.

Three extensions remain open:
formal grounding of independence, integrity, and controllability;
grounding of the capability model in claim resolution;
and comparison of evaluator contexts for an external framework.

% !TeX root = ../se320_structural_assurability.tex
% ==================================================
\section{Threats to Validity and Limitations}
\label{aa.sec.threats}
% ==================================================

\paragraph{Construct validity.}
The seven dimensions are provisional.
The separating-construction criterion requires a controlled
structural-profile difference and an independently established
capability difference.
A completed construction demonstrates only the claim-relative contribution,
not universal necessity or sufficiency.

\paragraph{Formal model.}
The capability model is abstract, and monotonicity is assumed.
Capability non-equivalence is a separate hypothesis from the profile difference.
Only four of the seven dimensions have dimension-specific constructions.
The Lean formalization checks the stated theorems.
It does not check that the models are faithful to a real system or claim.

\paragraph{External validity.}
The worked examples are finite and use table-lookup capabilities.
The pilot results are conditional on researcher-defined claims and admissibility assumptions.
No operational AI system has been assessed.

\paragraph{Exploratory experiments.}
The finite-model experiments are illustrations, not evidence about deployed systems.

% !TeX root = ../se320_structural_assurability.tex
% ==================================================
\section{Conclusion}
\label{aa.sec.conclusion}
% ==================================================

This paper proposed structural assurability:
the claim-relative capacity of a system and its evidence-producing environment
to make assurance-relevant evidence practically obtainable.
Assurability is compared by containment of capability sets, so incomparability is permitted.
The accompanying Lean formalization establishes the resolution bound
and separating constructions for four of the seven provisional dimensions.
Independence, integrity, controllability, and grounding of the capability model in claim resolution remain open.


% ==================================================
\section*{Statements and Declarations}
% ==================================================

\subsection*{Acknowledgments}
This research made use of SciX,
a scientific literature search and discovery platform,
whose related-literature tools helped identify relevant work across disciplinary boundaries.

\subsection*{Author Contributions}
The author is the sole contributor to this work
and is responsible for all aspects of the research, authorship, and publication.

\subsection*{Code Availability}
A reference implementation that checks the finite constructions of this paper
is available~\citep{case2026operationalidentityfiniteaudit}.

\subsection*{Use of AI-Assisted Tools}
AI-assisted tools were used for editing, formatting, and consistency checking,
and in developing and reviewing mathematical definitions, proofs,
and the companion Lean~4 formalization.
The author reviewed all suggestions and is solely responsible for the content.

\subsection*{Declaration of Conflicting Interest}
The author confirms there are no potential conflicts of interest.

% ==================================================
% References
% ==================================================

\bibliographystyle{plainnat}
\bibliography{references}

\end{document}
